\documentclass[11pt]{article}
\usepackage{fontspec}
% Charter et Menlo (macOS) ; ailleurs, XCharter et DejaVu Sans Mono, dont ils dérivent.
\IfFontExistsTF{Charter}{\setmainfont{Charter}}{\setmainfont{XCharter}}
\IfFontExistsTF{Menlo}{\setmonofont{Menlo}[Scale=0.82]}{\setmonofont{DejaVu Sans Mono}[Scale=0.82]}
\usepackage{polyglossia}\setmainlanguage{french}\setotherlanguage{english}\setotherlanguage{latin}
\usepackage{microtype}
\usepackage[most]{tcolorbox}
\usepackage{etoolbox}
\usepackage{titlesec}
\titleformat{\part}[display]{\normalfont\bfseries}{\large Partie~\arabic{part}}{0.4em}{\huge}
\titlespacing*{\part}{0pt}{0pt}{2em}
\pretocmd{\part}{\clearpage}{}{}
\newtcolorbox{encadre}{enhanced, breakable, colback=white, colframe=black!55, boxrule=0.6pt, arc=0pt, left=6pt, right=6pt, top=2pt, bottom=2pt, before skip=8pt, after skip=8pt}
\usepackage[a4paper,margin=27mm]{geometry}
\usepackage{booktabs,longtable,array}
\usepackage{setspace}\setstretch{1.1}
\usepackage{xcolor}
\usepackage[hidelinks]{hyperref}
\usepackage{amsmath,amssymb,amsthm}

\newcommand{\q}[1]{«\,#1\,»}
% Une vue de Gallica dans la vue de lecture du site : volume, lot, vue.
\newcommand{\vue}[3]{\href{https://germain.commutator.io/archive/\##1/#2/fr/#3}{#1, lot~#2, vue~#3}}
% Le fichier de la preuve vérifiée, une fois par section.
\newcommand{\leanblob}{https://github.com/Commutator-IO/germain-project/blob/main/lean/Germain}
\newcommand{\verifie}[1]{\footnote{Preuve vérifiée : \href{\leanblob/#1.lean}{\texttt{lean/Germain/\detokenize{#1}.lean}}.}}
% Une piste : un énoncé que l'onglet Findings du site tient pour possiblement absent
% de la littérature. Non établi comme nouveau ; le losange renvoie à la fiche.
\definecolor{piste}{HTML}{B53D1D}
\newcommand{\piste}[1]{\,\href{https://germain.commutator.io/findings/\##1}{\textcolor{piste}{\ensuremath{\blacklozenge}}}}
\definecolor{pisteresolue}{gray}{0.55}
\newcommand{\pisteresolue}[1]{\,\href{https://germain.commutator.io/findings/\##1}{\textcolor{pisteresolue}{\ensuremath{\blacklozenge}}}}
\theoremstyle{plain}
\newtheorem*{enonce}{Énoncé}
\newtheorem{lemme}{Lemme}[section]
\newtheorem{theoreme}[lemme]{Théorème}
\newtheorem{contrexemple}[lemme]{Contre-exemple}
\newtheorem{corollaire}[lemme]{Corollaire}
\theoremstyle{definition}
\newtheorem*{definition}{Définition}
\theoremstyle{remark}
\newtheorem*{remarque}{Remarque}
\raggedbottom\emergencystretch=1.5em

\title{\Large Énoncés des lectures modernisées démontrés en Lean\\[4pt]\large Manuscrits de mathématiciens dans Gallica}
\author{Michel Hua\\ \small Chercheur indépendant, Paris\\ \small \texttt{michel@commutator.io}}
\date{\small Octobre 2026}

\begin{document}
\maketitle

\begin{abstract}
\noindent Les lectures modernisées de germain.commutator.io reformulent en langage actuel les manuscrits de mathématiciens que Gallica numérise en entier. Elles dérivent d'une transcription faite par un modèle et peuvent se tromper. Ce document démontre des énoncés de ces lectures. Chaque démonstration suit une preuve vérifiée dans l'assistant de preuve Lean~4, avec sa bibliothèque mathlib. Pour chaque énoncé, le document donne le contexte du feuillet, la démonstration, les hypothèses superflues ou nécessaires, et les corrections apportées à la lecture ou au feuillet.
\end{abstract}

\tableofcontents

\section*{Introduction}\addcontentsline{toc}{section}{Introduction}

Une preuve vérifiée établit que l'énoncé tient sous les hypothèses écrites. Elle ne dit pas que la lecture est fidèle au feuillet, ce qui demande le fac-similé, ni que l'énoncé manque à la littérature. Le losange rouge \textcolor{piste}{\ensuremath{\blacklozenge}} marque les énoncés que l'onglet Findings du site retient comme pistes : des énoncés qui pourraient manquer à la littérature publiée. Une piste n'est ni un résultat établi ni une affirmation de priorité ; le losange renvoie à sa fiche. Chaque lien \vue{fr-9115}{1}{1} conduit à la vue de Gallica dans la vue de lecture du site, à côté de sa transcription.

Les parties suivent les mathématiciens dans l'ordre chronologique. Ce premier état ne contient que Sophie Germain ; les autres viendront.

\part{Sophie Germain (1776--1831)}
\section[Français 9115 : le manuscrit C et les résidus quadratiques]{Français 9115 : le manuscrit C et les résidus quadratiques\protect\verifie{Fr9115Arithmetic}}\label{sec:fr-9115-arithmetique}

Vers la fin du volume Français 9115, entre les vues 704 et 730, sont reliés trois pages d'une mise au net sur des équations de Fermat à exposants pairs, le « manuscrit C » de Laubenbacher et Pengelley (ff. 348r--349r, \vue{fr-9115}{36}{708} à \vue{fr-9115}{36}{710}), et, autour d'elles, des feuillets et des billets de travail sur les résidus quadratiques (\vue{fr-9115}{36}{704} à \vue{fr-9115}{37}{730}). La lecture modernisée restitue l'argument du manuscrit C pas à pas et relève quatre points où il ne tient pas ; elle vérifie ensuite, billet par billet, les énoncés sur les résidus. On démontre ici ce qui, dans ces deux parties de la lecture, se laisse énoncer et démontrer : les identités et les divisibilités sur lesquelles repose le manuscrit C, des contre-exemples aux pas qui échouent, et la plupart des énoncés des notes sur les résidus, avec leurs exceptions.

\noindent\emph{Conventions.} « $a$ est résidu de $p$ » se lit : $a$ est un carré dans $\mathbf Z/p\mathbf Z$ (\texttt{IsSquare}). Avec cette convention, $0$ est un carré : un multiple de $p$ est résidu de $p$. On note $\left(\frac{a}{n}\right)$ le symbole de Jacobi pour $n$ impair positif, qui est le symbole de Legendre quand $n$ est premier. Les nombres de Fermat sont $F_k = 2^{2^k}+1$.

\subsection*{Le manuscrit C}

\begin{contrexemple}[Sophie Germain, vue 708, l'énoncé tel qu'il est écrit]
Pour tout $n$, $x = y = z = 1$ satisfait $2z^{2n} = x^{2n} + y^{2n}$ avec $x$, $y$ premiers entre eux ; l'énoncé « il est impossible de satisfaire en nombres entiers à l'équation $2z^{2n} = y^{2n} + x^{2n}$ » est donc faux tel quel. Pour $n = 1$, on a en outre $1^2 + 7^2 = 2\cdot 5^2$ et $7^2 + 17^2 = 2 \cdot 13^2$.
\end{contrexemple}

C'est \texttt{manuscriptC\_not\_as\_written} et \texttt{manuscriptC\_n\_one}. L'énoncé corrigé (seules solutions primitives, $x^2 = y^2 = z^2 = 1$, pour $n$ impair $> 1$) est le théorème de Darmon et Merel ; il n'est pas démontré ici.

\begin{enonce}[Énoncé de l'édition, vue 708, les hypothèses]
Soit $n \geq 1$. Si $x$, $y$ sont premiers entre eux et $x^{2n} + y^{2n} = 2z^{2n}$, alors $x$, $y$ et $z$ sont impairs.
\end{enonce}

\begin{proof}
On réduit modulo $4$, où le carré d'un impair vaut $1$ et celui d'un pair $0$. Si $x$ et $y$ étaient pairs, $2$ diviserait leurs coefficients de Bézout, ce qui est exclu. Si exactement l'un est pair, le membre de gauche vaut $1$ modulo $4$, et le membre de droite $0$ ou $2$. Si $x$ et $y$ sont impairs et $z$ pair, le membre de gauche vaut $2$ et le membre de droite $0$ (\texttt{manuscriptC\_odd}).
\end{proof}

\begin{enonce}[Sophie Germain, vue 708, les paramétrisations et le système (1)]
On a les identités
\[ (P^2-Q^2-2PQ)^2 + (P^2-Q^2+2PQ)^2 = 2(P^2+Q^2)^2, \]
et, pour $x = p^2-q^2-2pq$, $y = p^2-q^2+2pq$ : $x + y = 2(p^2-q^2)$, $y - x = 4pq$, $y^2 - x^2 = 8\,pq(p^2-q^2)$.
\end{enonce}

Ce sont des identités polynomiales (\texttt{param\_identity}, \texttt{system\_one}). La dernière donne le $N = \tfrac18(y^2-x^2)$ de la lecture.

\begin{enonce}[Énoncé de l'édition, vue 708, la factorisation et la première lacune]
Pour $n$ impair, posons $\Phi = \sum_{i=0}^{n-1} (x^2)^i (-y^2)^{n-1-i} = x^{2n-2} - x^{2n-4}y^2 + \cdots + y^{2n-2}$. Alors $x^{2n} + y^{2n} = (x^2+y^2)\,\Phi$ et $\Phi \equiv n\,y^{2n-2} \pmod{x^2+y^2}$. Si $x$ et $y$ sont premiers entre eux, tout diviseur commun de $x^2+y^2$ et de $\Phi$ divise $n$.
\end{enonce}

\begin{proof}
La première égalité est l'identité $(\sum_i a^i b^{n-1-i})(a-b) = a^n - b^n$ appliquée à $a = x^2$, $b = -y^2$, avec $(-y^2)^n = -y^{2n}$ pour $n$ impair (\texttt{factor\_Phi}). Comme $n-1$ est pair, $y^{2n-2} = (-y^2)^{n-1}$, et
\[ \Phi - n\,y^{2n-2} = \sum_{i=0}^{n-1} \bigl((x^2)^i - (-y^2)^i\bigr)(-y^2)^{n-1-i}, \]
où chaque terme est divisible par $x^2 - (-y^2) = x^2+y^2$ (\texttt{Phi\_modEq}). Un diviseur commun $d$ divise donc $n\,y^{2n-2}$. De $ax + by = 1$ on tire $a^2(x^2+y^2) + (2b - a^2y - b^2y)\,y = 1$, de sorte que $x^2+y^2$, et donc $d$, est premier avec $y$ et avec $y^{2n-2}$ : $d$ divise $n$ (\texttt{dvd\_n\_of\_dvd\_both}).
\end{proof}

\begin{contrexemple}[Énoncé de l'édition, vue 708]
Pour $n = 5$, $x = 1$, $y = 3$ : $x^2+y^2 = 10$, $\Phi = 5905 = 5 \cdot 1181$, et $5$ divise les deux.
\end{contrexemple}

C'est \texttt{Phi\_counterexample}. Le partage $x^2+y^2 = 2z'^{2n}$, $\Phi = z''^{2n}$ demande que $x^2+y^2$ et $\Phi$ soient premiers entre eux au facteur $2$ près ; l'exemple montre que cela peut manquer quand un facteur premier de $n$ est de la forme $4k+1$. Ce n'est pas une solution de l'équation ; il montre seulement que le pas ne va pas de soi. Laubenbacher et Pengelley avaient déjà relevé cette lacune.

\begin{enonce}[Sophie Germain, vue 708, le système (2)\piste{fr-9115-manuscript-c-case-one}]
Soit $s = \pm 1$, $P = Uq + sVp$, $Q = Up - sVq$. Alors $P^2 + Q^2 = (U^2+V^2)(p^2+q^2)$,
\[ P^2 - Q^2 = \bigl[U(p+q) + sV(p-q)\bigr]\bigl[U(q-p) + sV(p+q)\bigr] = (p^2-q^2)(V^2-U^2) + 4s\,pq\,UV, \]
\[ PQ = (U^2-V^2)\,pq + s\,(p^2-q^2)\,UV. \]
Si de plus $p$ et $q$ sont premiers entre eux et de parités différentes, et si $P^2 - Q^2 = H(p^2-q^2)$ et $PQ = F\,pq$, alors $pq(p^2-q^2)$ divise $UV$.
\end{enonce}

\begin{proof}
Les trois identités se vérifient en développant, avec $s^2 = 1$ (\texttt{system\_two}). Pour la divisibilité (\texttt{N\_dvd\_UV}), $p^2 - q^2 = (p+q)(p-q)$ est impair, donc premier avec $4$ ; de $ap + bq = 1$ on tire des relations de Bézout entre $p^2-q^2$ et $p$, et entre $p^2-q^2$ et $q$ (ainsi $-b^2(p^2-q^2) + (b^2p + 2a - a^2p)\,p = 1$). Donc $p^2-q^2$ est premier avec $4pq$. La première équation donne $s \cdot 4pq\,UV = (p^2-q^2)\bigl(H - (V^2-U^2)\bigr)$, et, en multipliant par $s$, $p^2-q^2 \mid 4pq\,UV$, d'où $p^2-q^2 \mid UV$. La seconde donne de même $pq \mid (p^2-q^2)\,UV$, d'où $pq \mid UV$. Comme $pq$ et $p^2-q^2$ sont premiers entre eux, leur produit divise $UV$.
\end{proof}

\begin{remarque}
Les hypothèses de coprimalité et de parité sur $p$, $q$ sont celles que la lecture tire de $x$, $y$ premiers entre eux et impairs ; la page ne les écrit pas. Que $P$ et $Q$ s'écrivent $Uq \pm Vp$, $Up \mp Vq$ avec $U^2 + V^2 = z''^n$ est supposé ici, comme sur la page ; la lecture le justifie par les entiers de Gauss, ce qui n'est pas formalisé.
\end{remarque}

\begin{enonce}[Sophie Germain, vue 709, les deux signes\piste{fr-9115-manuscript-c-case-one}]
Soit $s = \pm 1$. On a $(U'u + sV'v)(U'v - sV'u) = (U'^2-V'^2)\,uv + s\,(v^2-u^2)\,U'V'$. Si $U'$ et $V'$ sont premiers entre eux et si $N$ divise \emph{à la fois} $(U'^2-V'^2)\,uv + (v^2-u^2)\,U'V'$ et $(U'^2-V'^2)\,uv - (v^2-u^2)\,U'V'$, alors $N$ divise $2(v^2-u^2)\,uv$.
\end{enonce}

\begin{proof}
L'identité se vérifie en développant (\texttt{UV\_identity}). La somme et la différence des deux valeurs donnent $N \mid 2(U'^2-V'^2)\,uv$ et $N \mid 2(v^2-u^2)\,U'V'$. Les relations de Bézout du paragraphe précédent, appliquées à $U'$, $V'$, montrent que $U'^2 - V'^2$ est premier avec $U'$ et avec $V'$, donc avec $U'V'$ : $\alpha(U'^2-V'^2) + \beta\,U'V' = 1$. Alors
\[ 2(v^2-u^2)\,uv = \alpha(v^2-u^2)\cdot 2(U'^2-V'^2)\,uv + \beta\,uv\cdot 2(v^2-u^2)\,U'V', \]
et les deux termes sont divisibles par $N$ (\texttt{two\_signs}).
\end{proof}

\begin{remarque}
La conclusion de la page est juste \emph{sous l'hypothèse} que les deux valeurs soient divisibles par $N$. Cette hypothèse est la deuxième lacune relevée par la lecture : le signe $s$ est fixé par la décomposition, une seule des deux expressions est $UV$, et seule celle-là est divisible par $N$ d'après l'énoncé précédent. La démonstration ci-dessus ne comble pas cette lacune ; elle établit seulement que le reste du pas est correct.
\end{remarque}

\begin{contrexemple}[Sophie Germain, vue 709, la chaîne d'inégalités\piste{fr-9115-manuscript-c-case-one}]
On a $2(y^2-x^2) - (x+y)^2 = (x+y)(y-3x)$. Pour $p = 3$, $q = 2$, la paramétrisation donne $x = -7$, $y = 17$, avec $x^2+y^2 = 2\cdot 13^2$ : $(x+y)^2 = 100 < 480 = 2(y^2-x^2)$ mais $(x+y)^2 < y^2 = 289$. Pour $x = 7$, $y = 17$ : $2(y^2-x^2) = 480 < 576 = (x+y)^2$.
\end{contrexemple}

\texttt{inequality\_identity} et \texttt{inequality\_counterexamples} les établissent par le calcul. L'inégalité $2(y^2-x^2) > (x+y)^2$ que la page affirme n'est donc vraie que si $(x+y)(y-3x) > 0$, et la comparaison avec $y^2$ demande $(x+y)^2 \geq y^2$. Les deux exemples sont des solutions de $x^2+y^2 = 2z^2$ et non de l'équation de degré $2n$ : ils montrent que la chaîne, telle qu'écrite, n'utilise rien qui exclue ces cas. Ils ne montrent pas qu'une autre borne serait impossible.

\begin{enonce}[Énoncé de l'édition, vue 709, le cas oublié]
Pour $n$ impair et $x = \pm 1$, on a $x^n = x$.
\end{enonce}

\texttt{trivial\_case} le vérifie. Les égalités $x^n = x$, $y^n = y$ auxquelles aboutit le second cas de l'alternative n'entraînent donc pas $n = 1$ : elles ont lieu pour $x^2 = y^2 = 1$ et tout $n$ impair.

\begin{enonce}[Sophie Germain, vue 710, le théorème de Fermat pour les exposants pairs]
Soit $n \geq 1$. Si $X^{2n} + Y^{2n} = Z^{2n}$, l'un de $X$, $Y$ est pair. L'équation $Z^{2n} = X^{2n} + Y^{2n}$ n'a pas de solution en entiers non nuls si $n$ est pair, ni si $3$ divise $n$.
\end{enonce}

\begin{proof}
Si $X$ et $Y$ étaient impairs, le membre de gauche vaudrait $2$ modulo $4$, et une puissance paire vaut $0$ ou $1$ (\texttt{fermat\_not\_both\_odd}). Si $n$ est pair, $4$ divise $2n$, et une solution donnerait une solution de $a^4 + b^4 = c^4$ en entiers non nuls, que le théorème de Fermat pour l'exposant $4$ (mathlib, \texttt{fermatLastTheoremFour}) exclut ; si $3$ divise $n$, de même avec l'exposant $3$ (\texttt{fermatLastTheoremThree}) (\texttt{fermat\_two\_n\_of\_even}, \texttt{fermat\_two\_n\_of\_three\_dvd}).
\end{proof}

\begin{remarque}
Pour $n$ pair, la réduction de la page au premier théorème ne vaut pas : $Z + Y$ ne divise pas $Z^n + Y^n$ en général, ainsi $4 \nmid 3^2 + 1^2$ (\texttt{even\_n\_no\_factor}). Le cas $n$ pair est couvert ici par l'exposant $4$, que la page ne mentionne pas. Le cas $n$ impair sans facteur $3$, que la page ramène au premier théorème, demande le théorème de Fermat pour un exposant premier $\geq 5$, absent de mathlib.
\end{remarque}

\subsection*{Résidus quadratiques}

\begin{enonce}[Sophie Germain, vue 704, la méthode]
Soit $p$ un nombre premier impair et $m$ un diviseur de $p-1$. Si $m \equiv 1 \pmod 4$, $m$ est résidu de $p$ ; si $m \equiv 3 \pmod 4$, $-m$ est résidu de $p$. En conséquence, $2$ est non-résidu de $p$ si $p \equiv 3 \pmod 8$, résidu si $p \equiv 7 \pmod 8$, et résidu si $p = 2^{2k+3}w + 1$ avec $w$ impair.
\end{enonce}

\begin{proof}
Par la loi de réciprocité pour le symbole de Jacobi, $\left(\frac{m}{p}\right) = \left(\frac{p}{m}\right)$ si $m \equiv 1 \pmod 4$ ; comme $p \equiv 1 \pmod m$, cela vaut $\left(\frac{1}{m}\right) = 1$ (\texttt{residue\_of\_dvd\_pred\_one}). Si $m \equiv 3 \pmod 4$, $\left(\frac{-m}{p}\right) = (-1)^{(p-1)/2}\,(-1)^{\frac{m-1}{2}\frac{p-1}{2}}\left(\frac{p}{m}\right)$, et l'exposant total $\frac{p-1}{2}\bigl(1 + \frac{m-1}{2}\bigr)$ est pair (\texttt{neg\_residue\_of\_dvd\_pred\_three}).

Si $p \equiv 3 \pmod 8$, $m = \frac{p-1}{2} \equiv 1 \pmod 4$ est résidu, $-1$ est non-résidu, et $2m \equiv -1$ : si $2$ était résidu, $-1$ le serait (\texttt{two\_nonresidue\_of\_three\_mod\_eight}). Si $p \equiv 7 \pmod 8$, $m \equiv 3 \pmod 4$, $-m = s^2$ et $2 \cdot (-m) \equiv 1$, d'où $2 = (1/s)^2$ (\texttt{two\_residue\_of\_seven\_mod\_eight}). Si $p = 2^{2k+3}w + 1$, $-1 = t^2$ (car $p \equiv 1 \pmod 4$), $w$ est résidu ($w$ ou $-w$ l'est, et $-1$ aussi), $w = s^2$, et $2\,(2^{k+1}s)^2 = 2^{2k+3}w = -1 = t^2$ : $2 = \bigl(t/(2^{k+1}s)\bigr)^2$ (\texttt{two\_residue\_of\_odd\_two\_power}).
\end{proof}

\begin{remarque}
La méthode suppose la loi de réciprocité, que la page emploie sans la nommer. Elle ne dit rien quand la puissance de $2$ qui divise exactement $p - 1$ a un exposant pair, en particulier pour $p \equiv 5 \pmod 8$ : c'est ce que dit la dernière ligne de la page, et l'énoncé n'en couvre pas davantage. Le second complément lui-même, $\left(\frac{2}{p}\right) = 1 \iff p \equiv \pm 1 \pmod 8$, est dans mathlib (\texttt{ZMod.exists\_sq\_eq\_two\_iff}, repris en \texttt{two\_residue\_iff}).
\end{remarque}

\begin{enonce}[Énoncé de l'édition, vue 706, les nombres premiers $4n+1$]
Pour deux nombres premiers $P \equiv p \equiv 1 \pmod 4$, $P$ est résidu de $p$ si et seulement si $p$ est résidu de $P$ ; si $P$ est résidu de $p$, $-P$ l'est aussi. On a, pour $s = \pm 1$,
\[ (a^2+b^2)(\pi^2+\rho^2+\sigma^2+\tau^2) = (a\pi + sb\rho)^2 + (b\pi - sa\rho)^2 + (a\sigma + sb\tau)^2 + (a\tau - sb\sigma)^2. \]
\end{enonce}

La première partie est le cas $p \equiv q \equiv 1 \pmod 4$ de la loi de réciprocité (mathlib ; \texttt{four\_n\_one\_reciprocity}) ; la deuxième vient de ce que $-1$ est résidu de $p$ (\texttt{neg\_residue\_of\_one\_mod\_four}) ; l'identité se vérifie en développant (\texttt{four\_square\_expansion}). La page cherche la première partie par l'égalité terme à terme des deux membres de l'identité, ce qui n'est pas une inférence valable ; la démonstration ci-dessus ne suit pas la page.

\begin{enonce}[Sophie Germain, vue 722, la proposition générale]
Tout entier $a$ qui n'est pas un carré (en particulier tout $-t^2$, $t \neq 0$) est non-résidu d'un nombre premier $p$ qui ne le divise pas. Réciproquement, un carré est résidu de tout nombre premier.
\end{enonce}

\begin{proof}
Il suffit de trouver $N$ impair avec $\left(\frac{a}{N}\right) = -1$ : la multiplicativité du symbole de Jacobi en bas fournit alors un facteur premier $p$ de $N$ avec $\left(\frac{a}{p}\right) = -1$, donc $a$ non-résidu de $p$ et $p \nmid a$ (\texttt{exists\_prime\_nonresidue}). On a $a \neq 0$ ; écrivons $|a| = b^2 s$ avec $s$ sans facteur carré (\texttt{exists\_jacobiSym\_eq\_neg\_one}).

Si un nombre premier impair $q$ divise $s$, on écrit $s = qs'$ avec $q \nmid s'$, et $b = q^e b'$ avec $q \nmid b'$. Soit $c$ un non-résidu de $q$. Par le théorème chinois, il existe $N$ avec $N \equiv c \pmod q$ et $N \equiv 1 \pmod{4s'b'}$. Alors $N \equiv 1 \pmod 4$, $N$ est premier avec $q$ et avec $b'$, donc avec $b$ ; $\left(\frac{\pm1}{N}\right) = 1$, $\left(\frac{b^2}{N}\right) = 1$ ; $\left(\frac{s'}{N}\right) = \left(\frac{s'}{1}\right) = 1$, le symbole $\left(\frac{s'}{\cdot}\right)$ ne dépendant que de la classe modulo $4s'$ ; et $\left(\frac{q}{N}\right) = \left(\frac{N}{q}\right) = \left(\frac{c}{q}\right) = -1$, par la réciprocité ($N \equiv 1 \pmod 4$). Le produit vaut $-1$.

Sinon, $s$ n'a que $2$ pour facteur premier, et, étant sans facteur carré, vaut $1$ ou $2$. Écrivons $b = 2^j(2t+1)$. Le cas $a = b^2$ est exclu. Pour $a = -b^2$ et $a = 2b^2$, on prend $N = 4(2t+1)^2 - 1$, qui est $\equiv 3 \pmod 8$ et premier avec $b$ : $\left(\frac{-1}{N}\right) = -1$ et $\left(\frac{2}{N}\right) = -1$. Pour $a = -2b^2$, on prend $N = 4(2t+1)^2 + 1 \equiv 5 \pmod 8$, et $\left(\frac{-2}{N}\right) = -1$. La réciproque est immédiate (\texttt{isSquare\_zmod\_of\_isSquare}).
\end{proof}

\begin{remarque}
La démonstration est celle qu'on donne aujourd'hui, par le symbole de Jacobi (postérieur aux feuillets) ; elle ne suit pas la réduction de la page aux nombres $4n+1$, $4n+2$, $4n+3$, qui admet le cas $4n+1$. Les pas de cette réduction que la page démontre sont l'objet de l'énoncé suivant.
\end{remarque}

\begin{enonce}[Sophie Germain, vues 720 à 723, les nombres $4n+2$ et $4n+3$]
Tout $m \equiv 2 \pmod 4$ est non-résidu d'un nombre premier $q \equiv 3 \pmod 4$. Soit $N \equiv 3 \pmod 4$, $N \geq 7$, premier ou non. Alors $-N$ est non-résidu d'un nombre premier $q \equiv 3 \pmod 4$ plus petit que $N$, et $N$ est non-résidu d'un nombre premier plus petit que $N$.
\end{enonce}

\begin{proof}
Pour $m$ : $m+1 \equiv 3 \pmod 4$, donc $\left(\frac{-1}{m+1}\right) = -1$, et $m+1$ a un facteur premier $q$ dont $-1$ est non-résidu, c'est-à-dire $q \equiv 3 \pmod 4$ ; comme $m \equiv -1 \pmod q$, $m$ est non-résidu de $q$ (\texttt{four\_n\_two\_nonresidue}). C'est l'argument de la page avec $2k+1 = 1$.

Pour $-N$ : $N - 4 \equiv 3 \pmod 4$ est positif ; il a un facteur premier $q \equiv 3 \pmod 4$, $q \leq N - 4 < N$, et $N \equiv 4$, $-N \equiv -1\cdot 2^2 \pmod q$ ; si $-N$ était un carré, $-1$ en serait un, $2$ étant inversible (\texttt{neg\_nonresidue\_smaller}). Pour $+N$ : si $N \equiv 3 \pmod 8$, $N \geq 11$ et $N - 8 \equiv 3 \pmod 8$, donc $\left(\frac{2}{N-8}\right) = -1$, et un facteur premier $q$ de $N - 8$ a $2$ pour non-résidu ; $N \equiv 2 \cdot 2^2 \pmod q$ conclut. Si $N \equiv 7 \pmod 8$, $N - 2 \equiv 5 \pmod 8$, un facteur premier $q$ de $N-2$ a $2$ pour non-résidu, et $N \equiv 2 \pmod q$ (\texttt{pos\_nonresidue\_smaller}).
\end{proof}

\begin{remarque}
Comme $q$ divise $N - 4$ (ou $N - 8$, ou $N - 2$) et non $4$, $8$, $2$, il ne divise pas $N$ : la précaution que la lecture demande pour $N$ composé est remplie d'elle-même par ce choix. Il faut $N \geq 7$ : pour $N = 3$, le seul nombre premier plus petit est $2$, dont $3$ et $-3$ sont résidus (\texttt{three\_exception}). Le choix $a = 1$ dans $p - 4a^2$, $p - 8a^2$, $p - 2a^2$ est permis par la remarque finale de la page, qui n'exige que $2a < \sqrt p$ ou $2a < \sqrt{p/2}$.
\end{remarque}

\begin{contrexemple}[Sophie Germain, vue 720, l'extension]
$-5$, $-17$ et $-21$ sont résidus de tous les nombres premiers $q \equiv 3 \pmod 4$ plus petits qu'eux. L'énoncé « tout nombre de la forme $4n+1$ pris négativement est non résidu de quelque nombre premier moindre que lui et de la forme $4n+3$ » est donc faux.
\end{contrexemple}

\texttt{false\_extension} l'établit par énumération : $-17 \equiv 1^2 \pmod 3$, $\equiv 2^2 \pmod 7$, $\equiv 4^2 \pmod{11}$ ; $-21 \equiv 0 \pmod 3$ et $\pmod 7$, $\equiv 1 \pmod{11}$, $\equiv 6^2 \pmod{19}$. $17$ est premier, de sorte que l'extension échoue même pour $p$ premier. Que ce soient les seules exceptions en dessous de $200\,000$, comme le dit la lecture, n'est pas démontré ici.

\begin{contrexemple}[Sophie Germain, vues 718 et 726]
Pour $p = 19 = 12 \cdot 1 + 7$, le nombre pair immédiatement supérieur à $\sqrt{19/3}$ est $a = 4$, et $3a^2 - p = 29 > p$ ; $19$ est pourtant non-résidu de $7$. Le nombre $3$ est résidu de $11$ ($5^2 \equiv 3$), et $11$ n'est pas de la forme $3m + b^2$ avec $m$, $b$ entiers.
\end{contrexemple}

C'est \texttt{nineteen\_exception} et \texttt{three\_eleven} ; la seconde partie se voit modulo $3$, où $b^2 \equiv 2$ est impossible. La première réfute l'affirmation « $3aa - p$ sera $<$ que $p$ » du billet 354, la seconde l'énoncé que vise l'induction du feuillet 360.

\begin{enonce}[Sophie Germain, vue 728]
Soit $p$ premier, $p \equiv 3 \pmod 4$. Si $n$ est non-résidu de $p$, $-n$ est résidu. Si $p \equiv 11 \pmod{12}$, $3$ est résidu et $-3$ non-résidu de $p$. Tout nombre premier qui divise $a^2 + a + 1$ a $-3$ pour résidu, puisque $4(a^2+a+1) = (2a+1)^2 + 3$ ; aucun nombre premier $12n+11$ ne divise donc $a^2+a+1$.
\end{enonce}

\begin{proof}
Par le critère d'Euler, $n^{(p-1)/2} = -1$ et, $\frac{p-1}{2}$ étant impair, $(-n)^{(p-1)/2} = (-1)(-1) = 1$ (\texttt{neg\_residue\_of\_nonresidue}). La page passe par une racine primitive ; le critère d'Euler en est une forme. Pour $p \equiv 11 \pmod{12}$, la réciprocité pour deux premiers $\equiv 3 \pmod 4$ donne $\left(\frac{3}{p}\right) = -\left(\frac{p}{3}\right) = -\left(\frac{2}{3}\right) = 1$, et $-3 = -1 \cdot 3$ avec $-1$ non-résidu (\texttt{twelve\_n\_eleven}). Enfin $a^2 + a + 1 \equiv 0$ donne $-3 \equiv (2a+1)^2$ (\texttt{neg\_three\_of\_dvd}, \texttt{cube\_identity}, \texttt{twelve\_n\_eleven\_not\_dvd}).
\end{proof}

La dernière partie confirme le constat de la page, qu'on ne peut « rien conclure pour ces nombres » de la formule $\frac{4(a^3-1)}{a-1} = (2a+1)^2 + 3$.

\begin{enonce}[Sophie Germain, vue 729\piste{fr-9115-fermat-primes-residues}]
Si $p = 12n + 5$ et $3n + 1$ sont premiers, $3$ est racine primitive de $p$.
\end{enonce}

\begin{proof}
On a $n \geq 1$, puisque $1$ n'est pas premier. Par la réciprocité ($p \equiv 1 \pmod 4$), $\left(\frac{3}{p}\right) = \left(\frac{p}{3}\right) = \left(\frac{2}{3}\right) = -1$, donc $3^{(p-1)/2} = 3^{2(3n+1)} = -1$ par le critère d'Euler. L'ordre de $3$ divise $p - 1 = 4(3n+1)$ ; pour qu'il soit $p-1$, il suffit que $3^{(p-1)/r} \neq 1$ pour chaque facteur premier $r$ de $p-1$, soit $r = 2$ et $r = 3n+1$. Pour $r = 2$, c'est $-1 \neq 1$. Pour $r = 3n+1$, $3^4 = 1$ donnerait $p \mid 80$, donc $p \leq 80$ et $n \leq 6$, et aucun des nombres $17, 29, 41, 53, 65, 77$ ne divise $80$ (\texttt{three\_primitive\_root}).
\end{proof}

\begin{remarque}
La page exclut l'ordre $4$ par $3^2 + 1 \equiv 0$, soit $p \mid 10$, sans écarter $3^2 \equiv 1$ ; la démonstration ci-dessus écarte $3^4 \equiv 1$ d'un coup. L'hypothèse « $3n + 1$ premier » sert à réduire les facteurs premiers de $p - 1$ à $2$ et $3n+1$.
\end{remarque}

\begin{enonce}[Sophie Germain, vue 730, les nombres premiers de Fermat\piste{fr-9115-fermat-primes-residues}]
Pour $j < k$, $F_k \equiv 2 \pmod{F_j}$. Si $F_j$ et $F_k$ sont premiers, $F_j$ est résidu de $F_k$ si et seulement si $2$ est résidu de $F_j$. En conséquence, pour $F_k$ premier, $3$ est non-résidu de $F_k$ si $k \geq 1$, $5$ si $k \geq 2$, et $7$ si $k \geq 1$ ; mais $F_j$ est résidu de $F_k$ dès que $2 \leq j < k$ (ainsi $17$ est résidu de $257$ et de $65537$, et $257$ de $65537$). Enfin tout non-résidu d'un nombre premier de Fermat en est racine primitive : $3$ et $7$ sont racines primitives de tout nombre premier de Fermat $F_k$, $k \geq 1$, et $5$ de tout $F_k$ premier, $k \geq 2$.
\end{enonce}

\begin{proof}
Le produit $F_0 F_1 \cdots F_{k-1}$ vaut $F_k - 2$ (mathlib), d'où $F_j \mid F_k - 2$, et $F_j > 2$ (\texttt{fermatNumber\_mod}). Pour $k \geq 1$, $F_k \equiv 1 \pmod 4$, et la réciprocité donne $\left(\frac{F_j}{F_k}\right) = \left(\frac{F_k}{F_j}\right) = \left(\frac{2}{F_j}\right)$ (\texttt{fermat\_residue\_iff}). Pour $F_0 = 3$ et $F_1 = 5$, $2$ est non-résidu (\texttt{three\_nonresidue\_fermat}, \texttt{five\_nonresidue\_fermat}) ; pour $j \geq 2$, $8 \mid 2^{2^j}$, $F_j \equiv 1 \pmod 8$ et $2$ est résidu (\texttt{fermat\_residue\_of\_two\_le}, \texttt{seventeen\_residue}, où $257$ et $65537$ sont prouvés premiers par calcul). Pour $7$ : $2^{2^k} \equiv 2$ ou $4 \pmod 7$ par récurrence (le carré échange $2$ et $4$), donc $F_k \equiv 3$ ou $5 \pmod 7$, non-résidus de $7$, et $\left(\frac{7}{F_k}\right) = \left(\frac{F_k}{7}\right) = -1$ (\texttt{seven\_nonresidue\_fermat}). Enfin, si $a$ est non-résidu de $F_k = 2 \cdot 2^m + 1$, $m = 2^k - 1$, on a $a^{2^m} = -1 \neq 1$ (critère d'Euler) et $a^{2^{m+1}} = 1$ ; l'ordre de $a$ est donc $2^{m+1} = F_k - 1$ (\texttt{fermat\_primitive\_root} et ses trois corollaires).
\end{proof}

\begin{remarque}
L'énoncé principal de la page, que les nombres premiers de Fermat plus petits sont non-résidus et donc racines primitives des plus grands, n'est vrai que pour $3$ et $5$ ; la correction interlinéaire « cela est faux » est confirmée par $17$, résidu de $257$. Le second paragraphe, $3 = 2^2 - 1$ et $7 = 2^3 - 1$ non-résidus de $F_i$, est vrai pour $i \geq 1$, bien que l'argument employé (exiger un carré exact) ait le défaut que la page vient de relever. La démonstration ci-dessus ne suit pas cet argument. Les énoncés portent sur les nombres de Fermat premiers : la primalité de $F_k$ est une hypothèse, et seule celle de $F_0$ à $F_4$ est établie.
\end{remarque}

Ne sont pas démontrés ici : le premier théorème du manuscrit C, corrigé de la solution triviale, qui est le théorème de Darmon et Merel, ni son extension à $2z^{2(n+2d)} = x^{2n} + y^{2n}$ ; le théorème de Fermat pour les exposants $2n$ avec $n$ impair sans facteur $3$, qui demande le théorème de Wiles ; la justification de $P = Uq \pm Vp$ et de $U = U'u \pm V'v$ par les entiers de Gauss ; le second cas de l'alternative ($v = 0$ ou $u = 0$) ; l'énoncé de Gauss (art.~125) que la page prête à « l'auteur », que tout nombre premier $p \equiv 1 \pmod 4$, $p > 5$, pris négativement est non-résidu d'un nombre premier plus petit ; les arguments des vues 718, 719, 722 et 724 sur les nombres $12n+7$, $12n+5$, $120k+1$, $3k-1$ et $24k+17$, hors l'exception $p = 19$ ; les vérifications numériques de la lecture (jusqu'à $5000$, $20\,000$ ou $200\,000$).

\section[Français 9115 : la lame vibrante, Euler et les « Remarques »]{Français 9115 : la lame vibrante, le mémoire d'Euler de 1779 et les « Remarques »\protect\verifie{Fr9115Elasticity}}\label{sec:fr-9115-elasticite}

Le volume Français 9115 contient, aux vues 144 à 206, une copie latine du mémoire d'Euler \emph{Investigatio motuum quibus laminae et virgae elasticae contremiscunt} (Acta Petropolitana, 1779), puis, aux vues 208 à 396, plusieurs états d'une critique de ses §§ 47 à 51, les « Remarques sur le mémoire d'Euler », et enfin, à partir de la vue 398, le mémoire sur la question proposée par la première classe de l'Institut, qui commence par l'équation de la plaque vibrante. On démontre ici les énoncés mathématiques que la lecture modernisée tire de ces feuillets : l'intégrale de l'équation du quatrième ordre, les équations aux fréquences, la position de quelques racines, les transformations de l'équation finale du § 47, le verdict sur la critique, l'équation de la lame libre touchée par un stylet, et les premières intégrales de la plaque.

\noindent\emph{Conventions.} L'abscisse réduite est $u = s/a \in [0,1]$, et l'on écrit
\[
Y(u) = a\cosh\omega u + b\sinh\omega u + c\cos\omega u + d\sin\omega u
\]
(\texttt{modeY}) ; c'est la forme d'Euler $\alpha e^{\omega u} + \beta e^{-\omega u} + \gamma\sin\omega u + \delta\cos\omega u$ avec $a = \alpha+\beta$, $b = \alpha-\beta$, $c = \delta$, $d = \gamma$ (\texttt{eulerY\_eq\_modeY}). Un bout est libre si $Y'' = Y''' = 0$, appuyé si $Y = Y'' = 0$, encastré si $Y = Y' = 0$. La seconde portion d'une lame coupée par un stylet en $u = \lambda$ s'écrit depuis l'autre bout, $Y_2(u) = \tilde Y(1-u)$ avec $\tilde Y$ de la même forme (\texttt{reflY}). On pose $\mu = 1-\lambda$, $P(t) = \cosh t\sin t - \sinh t\cos t$ et $C_+(t) = 1 + \cosh t\cos t$.

\subsection*{L'intégrale de l'équation et les six cas}

\begin{enonce}[Euler, § 9, \vue{fr-9115}{8}{153}]
$Y'''' = \omega^4 Y$. Si $f^4 = bc^2k$, $k > 0$ et $g \geq 0$, la vibration régulière $y = C\sin\bigl(\zeta + t\sqrt{2g/k}\bigr)\,Y(s/f)$ vérifie $\partial_t^2 y = -2g\,bc^2\,\partial_s^4 y$.
\end{enonce}

\begin{proof}
La dérivée de $Y$ est la fonction de même forme de coefficients $\omega(b, a, d, -c)$ : $\cosh$ et $\sinh$ s'échangent, $\cos$ et $\sin$ tournent. Appliquée quatre fois, cette substitution donne $\omega^4(a,b,c,d)$ (\texttt{iteratedDeriv\_four\_modeY}, \texttt{modeY\_ode}). Pour la vibration, on écrit le facteur du temps $C\sin(\zeta + \nu t) = C\sin\zeta\cos\nu t + C\cos\zeta\sin\nu t$, qui est de la même forme en $t$ avec $\omega = \nu$ : sa dérivée seconde est $-\nu^2$ fois lui-même, et $\nu^2 = 2g/k$. En $s$, la dérivée quatrième est $f^{-4}$ fois la fonction. Il reste $-\frac{2g}{k} = -2g\,bc^2 \cdot \frac{1}{bc^2k}$, ce que donne $f^4 = bc^2k$ (\texttt{regular\_vibration}).
\end{proof}

\begin{remarque}
Le facteur du temps est $\sin\bigl(\zeta + t\sqrt{2g/k}\bigr)$, comme la lecture le corrige ; la copie écrit le radical sur $2g$ seul. Avec $t\sqrt{2g}/k$, la dérivée seconde serait $-2g/k^2$ fois la fonction, et l'équation n'aurait lieu que pour $k = 1$.
\end{remarque}

\begin{enonce}[Euler, §§ 13 et 14, \vue{fr-9115}{8}{158} à \vue{fr-9115}{10}{190}]
Pour $\omega \neq 0$, il existe des coefficients $(a,b,c,d)$ non tous nuls satisfaisant les quatre conditions des deux bouts si et seulement si : $\cos\omega\cosh\omega = 1$ (cas I, libre et libre ; cas VI, encastré et encastré) ; $\sin\omega\cosh\omega = \cos\omega\sinh\omega$, c'est-à-dire $\tan\omega = \tanh\omega$ (cas II, libre et appuyé ; cas V, appuyé et encastré) ; $\cos\omega\cosh\omega = -1$ (cas III, libre et encastré) ; $\sin\omega = 0$ (cas IV, appuyé et appuyé).
\end{enonce}

\begin{proof}
La dérivée $k$-ième de $Y$ au point $u$ vaut $\omega^k$ fois le produit scalaire de $(a,b,c,d)$ avec une ligne fixe, $(\cosh x, \sinh x, \cos x, \sin x)$ pour $k = 0$ et ses images par la substitution précédente pour $k = 1, 2, 3$, en $x = \omega u$ (\texttt{iteratedDeriv\_modeV}). Comme $\omega \neq 0$, les quatre conditions forment un système linéaire homogène de matrice $4 \times 4$ (\texttt{bc\_iff}), qui a une solution non nulle si et seulement si son déterminant est nul (\texttt{exists\_mode\_iff}). Le calcul des déterminants (\texttt{det\_free\_free} et suivants) donne $2 - 2\cos\omega\cosh\omega$, $2(\cos\omega\sinh\omega - \cosh\omega\sin\omega)$, $2 + 2\cos\omega\cosh\omega$, $-4\sinh\omega\sin\omega$, $2(\sinh\omega\cos\omega - \sin\omega\cosh\omega)$ et $2 - 2\cos\omega\cosh\omega$. Pour le cas I, par exemple, on ajoute la première colonne à la troisième et la deuxième à la quatrième : le déterminant devient $(\cosh\omega - \cos\omega)^2 - (\sinh\omega - \sin\omega)(\sinh\omega + \sin\omega)$, qui vaut $2 - 2\cos\omega\cosh\omega$ par $\cosh^2 - \sinh^2 = 1$ et $\sin^2 + \cos^2 = 1$. Pour le cas IV, $\sinh\omega \neq 0$ (\texttt{case\_I} à \texttt{case\_VI}).
\end{proof}

\begin{enonce}[Euler, §§ 16, 17 et 30, \vue{fr-9115}{8}{160}, \vue{fr-9115}{9}{176}]
On a $\cos x\cosh x < 1$ pour $0 < x \leq 3\pi/2$. L'équation $\cos\omega\cosh\omega = 1$ a une seule racine dans $]0, 2\pi[$, et elle s'écrit $3\pi/2 + \varphi$ avec $0 < \varphi < 1/50$. L'équation $\cos\omega\cosh\omega = -1$ a une racine dans $]\pi/2, \pi[$.
\end{enonce}

\begin{proof}
La fonction $g(x) = \sin x\cosh x - \cos x\sinh x$ a pour dérivée $2\sin x\sinh x > 0$ sur $]0, \pi[$ et s'annule en $0$ : elle est positive sur $]0, \pi]$ (\texttt{P\_pos}). La dérivée de $f(x) = \cos x\cosh x$ est $-g(x)$, donc $f$ décroît strictement sur $[0, \pi]$ à partir de $f(0) = 1$ ; sur $[\pi, 3\pi/2]$, $\cos x \leq 0$ et $f \leq 0$ (\texttt{cos\_mul\_cosh\_lt\_one}). Ensuite $f(3\pi/2) = 0$ et $f(3\pi/2 + 1/50) = \sin\frac1{50}\cosh\bigl(\frac{3\pi}2 + \frac1{50}\bigr)$, où $\sin\frac1{50} > 0{,}0199$ et, avec $\pi > 3{,}141592$, $e > 2{,}7182818283$ et $e^{0{,}73} \geq 1 + 0{,}73 + 0{,}73^2/2$, le cosinus hyperbolique dépasse $54$ : le produit dépasse $1$, et le théorème des valeurs intermédiaires donne une racine dans l'intervalle. Sur $]3\pi/2, 2\pi[$, $\cos x > 0$ et $\sin x < 0$, donc $f' = \cos x\sinh x - \sin x\cosh x > 0$ et la racine y est unique ; avec la première partie, c'est la seule de $]0, 2\pi[$ (\texttt{first\_root\_free\_free}). Enfin $\cos x\cosh x + 1$ vaut $1$ en $\pi/2$ et $1 - \cosh\pi < 0$ en $\pi$ (\texttt{root\_free\_clamped}).
\end{proof}

\begin{remarque}
Euler fonde l'absence de racine dans le premier quadrant sur une série tronquée, ce qui n'est pas une preuve ; l'argument de monotonie la remplace et va jusqu'à $3\pi/2$. La borne $\varphi < 1/50$ (soit $1^\circ 09'$) s'accorde avec son $\varphi \approx 1/56$ mais ne décide pas de son « à peine un degré » : la racine calculée par la lecture, $3\pi/2 + 1^\circ 01'$, n'est pas vérifiée ici. De même, la valeur $107^\circ 26'$ de la racine du cas III et l'écart avec les $106^\circ 45'$ de la copie ne sont pas démontrés ; seul le quadrant l'est.
\end{remarque}

\subsection*{La lame sur trois appuis (§§ 47 à 51)}

\begin{enonce}[Euler, § 47, \vue{fr-9115}{10}{195} ; « Remarques », § 2, \vue{fr-9115}{17}{322}]
Soit $E(\lambda,\omega) = 2 - 2e^{2\omega} - \bigl(e^{\lambda\omega} - e^{(2-\lambda)\omega}\bigr)\bigl(e^{\lambda\omega} - e^{-\lambda\omega}\bigr)\frac{\sin\omega}{\sin\lambda\omega\,\sin\mu\omega}$ le second membre de l'équation finale d'Euler. Si $\sin\lambda\omega\,\sin\mu\omega \neq 0$,
\[
E(\lambda,\omega) = -\frac{4e^\omega\bigl(\sinh\omega\,\sin\lambda\omega\,\sin\mu\omega - \sin\omega\,\sinh\lambda\omega\,\sinh\mu\omega\bigr)}{\sin\lambda\omega\,\sin\mu\omega}.
\]
La forme $2(1 - e^{2\omega})\sin\lambda\omega\sin\mu\omega - (1 - e^{2\mu\omega})(e^{2\lambda\omega} - 1)\sin\omega$ de la mise au net est $E(\lambda,\omega)\sin\lambda\omega\sin\mu\omega$ ; pour $n\lambda = 1$, sa division par $e^{2\lambda\omega} - 1$ donne au signe près $2\bigl(1 + e^{2\lambda\omega} + \dots + e^{2(n-1)\lambda\omega}\bigr)\sin\lambda\omega\sin\mu\omega + (1 - e^{2\mu\omega})\sin\omega$. Si de plus $\sinh\lambda\omega\,\sinh\mu\omega \neq 0$, l'équation équivaut à $\coth\lambda\omega + \coth\mu\omega = \cot\lambda\omega + \cot\mu\omega$. Pour $\lambda = \frac12$ et $x = \omega/2$, $\sinh 2x\sin^2 x - \sin 2x\sinh^2 x = 2\sin x\sinh x\,(\cosh x\sin x - \cos x\sinh x)$, et le crochet du § 48 vaut $(1 + e^{2x})\sin x - (e^{2x} - 1)\cos x = 2e^x(\cosh x\sin x - \sinh x\cos x)$.
\end{enonce}

\begin{proof}
On pose $A = e^{\lambda\omega}$ et $B = e^{\mu\omega}$, de sorte que $e^\omega = AB$, $e^{2\omega} = A^2B^2$, $e^{(2-\lambda)\omega} = AB^2$, $e^{2\lambda\omega} = A^2$ et $e^{2\mu\omega} = B^2$ (\texttt{exp\_split}), et l'on exprime les sinus hyperboliques par $A$, $A^{-1}$, $B$, $B^{-1}$. Après multiplication par $\sin\lambda\omega\sin\mu\omega$, les deux membres de la première égalité sont des polynômes en ces quantités dont la différence est un multiple de $AA^{-1} - 1$ et de $BB^{-1} - 1$ (\texttt{eulerFinal\_mul}, \texttt{eulerFinal\_eq}). La deuxième tient parce que $(A - AB^2)(A - A^{-1}) = (1 - B^2)(A^2 - 1)$ dès que $AA^{-1} = 1$ (\texttt{germainForm\_eq}). La troisième est la somme géométrique $\bigl(\sum_{i<n} q^i\bigr)(q - 1) = q^n - 1$ avec $q = e^{2\lambda\omega}$ et $q^n = e^{2\omega}$ (\texttt{germain\_geom\_form}). Pour la quatrième, on développe $\sinh\omega = \sinh(\lambda\omega + \mu\omega)$ et $\sin\omega$ de même, et l'on réduit les deux sommes de quotients au même dénominateur (\texttt{final\_iff\_coth}). Les deux dernières sont les formules de duplication et $e^{2x} = e^x e^x$ (\texttt{half\_factor}, \texttt{euler\_half\_bracket}).
\end{proof}

\begin{enonce}[Sophie Germain, « Remarques », § 1, \vue{fr-9115}{11}{214}, \vue{fr-9115}{17}{322}\piste{fr-9115-euler-47-critique}]
Si $\tan\frac\omega2 = \tanh\frac\omega2$, écrit $\sin x\cosh x = \cos x\sinh x$ avec $x = \omega/2$, alors $\sin\omega\cosh\omega = \sinh\omega$, c'est-à-dire $\sin\omega = \tanh\omega$ ; et $\sin\omega = 0$ entraîne alors $\omega = 0$.
\end{enonce}

\begin{proof}
Par duplication, il s'agit de $2\sin x\cos x(\cosh^2x + \sinh^2x) = 2\sinh x\cosh x$. La différence des deux membres est $2(\cos x\cosh x - \sin x\sinh x)(\sin x\cosh x - \cos x\sinh x) + 2\sinh x\cosh x(\sin^2x + \cos^2x - 1)$, nulle sous l'hypothèse (\texttt{objection\_sin\_eq\_tanh}). Si $\sin\omega = 0$, il vient $\sinh\omega = 0$, donc $\omega = 0$ (\texttt{objection\_sin\_zero}).
\end{proof}

Le calcul de la critique est juste. Sa conclusion, que la seconde solution d'Euler « sort de l'état de la question », ne l'est pas. On note \emph{conditions du § 47} les huit conditions d'Euler pour la lame appuyée à ses deux bouts et touchée par un stylet en $\lambda$ : bouts appuyés, $y = 0$ des deux côtés du stylet, pente et courbure continues. La dérivée troisième n'y est pas supposée continue : son saut est la force du stylet, l'effort $F$ du § 8.

\begin{contrexemple}[à la conclusion du § 1 des « Remarques »]
Soit $\omega > 0$ une racine de $\tan\frac\omega2 = \tanh\frac\omega2$, et soient, pour $\lambda = \frac12$,
\[
Y_1(u) = \frac{\sinh\omega u}{\sinh\frac\omega2} - \frac{\sin\omega u}{\sin\frac\omega2}, \qquad Y_2(u) = Y_1(1-u).
\]
Alors $Y_1$ et $Y_2$ satisfont les conditions du § 47, $Y_1$ n'est pas identiquement nulle, et $Y_1'''(\frac12) \neq Y_2'''(\frac12)$ : c'est un vrai mode de la lame tenue en son milieu, et le stylet y exerce une force.
\end{contrexemple}

\begin{proof}
D'abord $\sinh\frac\omega2 > 0$, et $\sin\frac\omega2 \neq 0$ : sinon l'hypothèse donnerait $\cos\frac\omega2\sinh\frac\omega2 = 0$, donc $\cos\frac\omega2 = 0$, contre $\sin^2 + \cos^2 = 1$. L'hypothèse s'écrit $\coth\frac\omega2 = \cot\frac\omega2$. En $0$, $\sinh$ et $\sin$ s'annulent avec leurs dérivées secondes : $Y_1(0) = Y_1''(0) = 0$, et de même $Y_2(1) = Y_2''(1) = 0$. En $\frac12$, $Y_1 = 1 - 1 = 0$ et $Y_2 = 0$. La dérivée de $Y_2$ en $\frac12$ est l'opposée de celle de $Y_1$, qui vaut $\omega\bigl(\coth\frac\omega2 - \cot\frac\omega2\bigr) = 0$ : les pentes sont égales, toutes deux nulles. Les dérivées secondes, égales par symétrie, valent $2\omega^2$. Les dérivées troisièmes sont opposées, et $Y_1'''(\frac12) = \omega^3\bigl(\coth\frac\omega2 + \cot\frac\omega2\bigr) = 2\omega^3\coth\frac\omega2 > 0$. Enfin le coefficient de $\sinh\omega u$ est non nul, et les quatre fonctions sont indépendantes pour $\omega \neq 0$ (\texttt{modeY\_eq\_zero\_iff}), donc $Y_1 \neq 0$ (\texttt{second\_family\_is\_mode}).
\end{proof}

\begin{remarque}
Les deux fonctions, de la forme d'Euler, sont les modes symétriques de la lecture : chaque moitié vibre comme une lame appuyée à un bout et encastrée au stylet, ainsi que l'écrit le § 51. Le contre-exemple établit le verdict de la lecture sur la seconde solution ; il ne dit rien de l'hypothèse « purement exclutif » de la vue 214, qui est une lecture du feuillet.
\end{remarque}

\begin{enonce}[Énoncé de l'édition, d'après Euler, § 50, \vue{fr-9115}{10}{200} et \vue{fr-9115}{11}{201}\piste{fr-9115-euler-1779-printed-slips}]
Pour $\omega = 2i\pi$, $i$ entier non nul, $Y_1 = Y_2 = \sin\omega u$ (soit $\gamma = \gamma' = 1$, les autres coefficients nuls) satisfait les conditions du § 47 avec $y'''$ continu ; avec $Y_2 = -\sin\omega u$ sur la même abscisse ($\gamma' = -\gamma$, le signe imprimé), les pentes en $\frac12$ diffèrent.
\end{enonce}

\begin{proof}
$\sin 2i\pi = \sin i\pi = 0$ donne les conditions aux bouts et au stylet ; les autres sont des égalités d'une fonction avec elle-même. Les pentes en $\frac12$ sont $\pm\omega\cos i\pi$, et $\cos i\pi \neq 0$ puisque $\sin i\pi = 0$ (\texttt{first\_family\_sign}).
\end{proof}

\begin{enonce}[Sophie Germain, « Remarques », §§ 2 et 5, \vue{fr-9115}{17}{324} et \vue{fr-9115}{17}{330}]
Il existe un entier $k \geq 1$ tel que $\sin k\pi\lambda = 0$ si et seulement si $\lambda$ est rationnel : les modes $\sin k\pi u$ de la lame sans stylet qui ont un nœud au stylet, ses points « sympathiques », n'existent qu'aux positions rationnelles, et à toutes. D'autre part, pour un entier $m$, il existe des entiers $k$ et $j$ avec $(4k+1)m = 4j+1$ si et seulement si $m \equiv 1 \pmod 4$.
\end{enonce}

\begin{proof}
$\sin k\pi\lambda = 0$ équivaut à $k\pi\lambda = n\pi$ avec $n$ entier, d'où $\lambda = n/k$ ; réciproquement, si $\lambda = p/q$ on prend $k = q$ (\texttt{sympathetic\_iff\_rational}). Pour la seconde partie, l'égalité donne $m = 4(j - km) + 1$ ; réciproquement $k = 0$ convient (\texttt{common\_term\_iff}).
\end{proof}

\begin{remarque}
La seconde partie est l'arithmétique du § 5 : avec $\lambda = 1/\lambda'$ et $m = \lambda' - 1$, les deux suites de sons approchées, en $(4k+1)^2/\lambda^2$ et $(4j+1)^2/\mu^2$, ont un terme commun si et seulement si $m \equiv 1 \pmod 4$, et aucun quand $m$ est pair ou de la forme $4n+3$, comme le dit le feuillet ; pour $\lambda' = 6$ du brouillon, $m = 5$. Cela porte sur les racines approchées $(4k+1)\pi/4$ d'Euler, non sur les racines exactes.
\end{remarque}

\subsection*{La lame libre touchée par un stylet (§§ 8 et 9)}

On note \emph{conditions du § 8} les huit conditions de la lame libre à ses deux bouts et touchée par un stylet en $\lambda$ : $Y_1'' = Y_1''' = 0$ en $0$, $Y_2'' = Y_2''' = 0$ en $1$, $Y_1 = Y_2 = 0$ en $\lambda$, pente et courbure continues en $\lambda$.

\begin{enonce}[Sophie Germain, « Remarques », § 8, \vue{fr-9115}{18}{342}\piste{fr-9115-free-lamina-stylet-equation}]
Soient $\omega > 0$ et $0 < \lambda < 1$. Il existe un mouvement non identiquement nul satisfaisant les conditions du § 8 si et seulement si
\[
C_+(\mu\omega)\,P(\lambda\omega) + C_+(\lambda\omega)\,P(\mu\omega) = 0 .
\]
\end{enonce}

\begin{proof}
On écrit $Y_1$ avec les coefficients $(a,b,c,d)$ et $Y_2(u) = \tilde Y(1-u)$ avec $(a',b',c',d')$. Les bouts libres donnent $c = a$, $d = b$, $c' = a'$, $d' = b'$. Posons, en $t = \lambda\omega$, $X = \cosh t + \cos t$, qui est $> 0$ pour $t \neq 0$, et $Z = \sinh t + \sin t$, et de même $X'$, $Z'$ en $t' = \mu\omega$. La condition $Y_1(\lambda) = 0$ s'écrit $aX + bZ = 0$. Au moyen de cette relation et des identités $\cosh^2 - \sinh^2 = 1$, $\sin^2 + \cos^2 = 1$, on trouve $X\,Y_1'(\lambda) = 2\omega b\,C_+(t)$ et $X\,Y_1''(\lambda) = -2\omega^2 b\,P(t)$, et de même $X'\,Y_2'(\lambda) = -2\omega b'\,C_+(t')$ et $X'\,Y_2''(\lambda) = -2\omega^2 b'\,P(t')$. La continuité de la pente et celle de la courbure deviennent
\[
b\,C_+(t)\,X' + b'\,C_+(t')\,X = 0, \qquad b\,P(t)\,X' - b'\,P(t')\,X = 0 .
\]
Soit $G = C_+(t')P(t) + C_+(t)P(t')$. Ces deux relations donnent $bX'G = 0$ et $b'XG = 0$. Si $b \neq 0$ ou $b' \neq 0$, $G = 0$. Si $b = b' = 0$, alors $aX = 0$ et $a'X' = 0$, et tous les coefficients sont nuls ; or un mouvement non nul a des coefficients non tous nuls. Réciproquement, si $G = 0$, on cherche $b = BX$, $a = -BZ$, $b' = B'X'$, $a' = -B'Z'$ avec $B\,C_+(t) + B'\,C_+(t') = 0$ et $B\,P(t) - B'\,P(t') = 0$. Si $C_+(t') \neq 0$, le couple $(B, B') = (C_+(t'), -C_+(t))$ convient ; sinon $P(t') \neq 0$, parce que $C_+$ et $P$ ne s'annulent pas ensemble hors de $0$ (\texttt{not\_both}), et $(P(t'), P(t))$ convient. On a $b \neq 0$, et $Y_1 \neq 0$ par l'indépendance des quatre fonctions (\texttt{core\_forward}, \texttt{core\_backward}, \texttt{free\_stylet\_iff}).
\end{proof}

\begin{remarque}
C'est l'équation du feuillet, qui l'écrit avec les exponentielles comme somme de deux produits, le second obtenu en échangeant $\lambda$ et $1-\lambda$ ; ses accolades valent quatre fois les facteurs $C_+$ et $P$, ce qui exige le terme de tête 4 que la transcription marque comme surchargé ou douteux. La preuve élimine d'abord, comme le feuillet, les conditions des bouts puis $Y(\lambda) = 0$, sans passer par le déterminant $8 \times 8$, qui n'est pas formalisé ; elle ne vérifie pas les deux valeurs de $\alpha/\alpha'$ du feuillet.
\end{remarque}

\begin{contrexemple}[à la conclusion du § 9, \vue{fr-9115}{18}{344}]
Pour $\lambda = \frac12$ et toute racine $\omega > 0$ de $C_+(\omega/2) = 0$, il existe un mouvement non nul satisfaisant les conditions du § 8, que le feuillet rejette. Ce que le feuillet démontre est juste : pour $x \neq 0$, $C_+(x) = 0$ et $\cos 2x\cosh 2x = 1$ sont incompatibles.
\end{contrexemple}

\begin{proof}
Pour $\lambda = \mu = \frac12$, l'équation devient $2\,C_+(\omega/2)P(\omega/2) = 0$, satisfaite (\texttt{free\_stylet\_half}). Pour l'incompatibilité, avec $c = \cosh x$, $\cos 2x = 2\cos^2x - 1$ et $\cosh 2x = \cosh^2x + \sinh^2x$, et $c\cos x = -1$ ; l'égalité $\cos 2x\cosh 2x = 1$ entraîne alors $(c^2 - 1)^2 = 0$, donc $\sinh x = 0$ et $x = 0$ (\texttt{cplus\_incompatible}).
\end{proof}

\begin{remarque}
Les mouvements de la première famille de son § 9, $\tan\frac\omega2 = \tanh\frac\omega2$, sont aussi donnés par l'énoncé du § 8 ; ce sont des mouvements de la lame libre, ce que la lecture vérifie et qui n'est pas démontré ici. Ceux de la seconde, $C_+(\omega/2) = 0$, ne le sont pas, comme le feuillet le montre ; mais ce sont des mouvements de la lame touchée par le stylet, et le feuillet n'avait pas à les rejeter.
\end{remarque}

\subsection*{La plaque : les §§ 1 et 2 du mémoire pour l'Institut}

On note $\Delta = \partial_x^2 + \partial_y^2$ ; l'équation (B) des mouvements réguliers s'écrit $Kb\tau^2\Delta^2 z = z$, ou $h^4\Delta^2 z = z$ avec $h^4 = Kb\tau^2$.

\begin{enonce}[Sophie Germain, N\textsuperscript{o} 2, \vue{fr-9115}{21}{404}]
Pour tout opérateur linéaire $L$ (le laplacien sur les fonctions régulières), une solution de $h^2Lz = z$ ou de $h^2Lz = -z$ est solution de $h^4L^2z = z$ ; réciproquement, toute solution de $h^4L^2z = z$ est la somme de $z' = \frac12(z + h^2Lz)$, solution de la première, et de $z'' = \frac12(z - h^2Lz)$, solution de la seconde.
\end{enonce}

\begin{proof}
$h^4L^2z = h^2L(h^2Lz)$, et l'on remplace deux fois $h^2Lz$ par $\pm z$ (\texttt{factor\_B\_left}). Pour la réciproque, $h^2Lz' = \frac12(h^2Lz + h^4L^2z) = \frac12(h^2Lz + z) = z'$, et de même $h^2Lz'' = -z''$ (\texttt{factor\_B\_split}).
\end{proof}

\begin{enonce}[Sophie Germain, N\textsuperscript{os} 3 et 4, intégrales (C) et (D), \vue{fr-9115}{21}{405} et \vue{fr-9115}{21}{406}]
Si $f$ et $g$ sont de la forme d'Euler en la variable $kx$ et $ky$, $z = f(x) + g(y)$ vérifie $\Delta^2 z = k^4 z$, donc (B) si $Kb\tau^2k^4 = 1$. Pour $z = \sin\alpha x\,\sin\beta y$, $\Delta z = -(\alpha^2 + \beta^2)z$ et $\Delta^2 z = (\alpha^2 + \beta^2)^2z$. Avec $\alpha = \pi n/A$, $\beta = \pi m/A$ et $n$, $m$ entiers, $z$ s'annule sur les quatre côtés du carré $[0, A]^2$, et si $Kb\tau^2(\alpha^2 + \beta^2)^2 = 1$ et $\tau > 0$, son son vaut $\frac{\sqrt{1/K}}{\pi} = \frac{\pi(m^2 + n^2)\tau\sqrt b}{A^2}$.
\end{enonce}

\begin{proof}
Pour une fonction de $x$ plus une fonction de $y$, le laplacien est la somme des dérivées secondes (\texttt{lap\_sep}) ; pour un produit, $\Delta(FG) = F''G + FG''$ (\texttt{lap\_prod}). Appliqués deux fois avec la règle de dérivation de la forme d'Euler, ils donnent $k^4z$ (\texttt{plate\_C}) et $(\alpha^2+\beta^2)^2z$, la fonction $\Delta z$ étant encore un produit de sinus (\texttt{plate\_D}). Sur les côtés, $\sin 0 = 0$ et $\sin(\pi n) = 0$ (\texttt{plate\_D\_boundary}). Enfin $1/K = \bigl(\tau(\alpha^2+\beta^2)\bigr)^2 b$, d'où $\sqrt{1/K} = \tau(\alpha^2+\beta^2)\sqrt b$ et $\alpha^2 + \beta^2 = \pi^2(n^2+m^2)/A^2$ (\texttt{plate\_D\_sound}).
\end{proof}

\begin{remarque}
Par \texttt{plate\_D}, le laplacien de (D) s'annule aussi sur les côtés : (D) est un mode de la plaque carrée appuyée sur son contour, quelle que soit la matière, comme le dit la lecture. Dans (C), on a pris deux fonctions d'Euler quelconques de même $k$ ; le feuillet prend la même fonction $f$ multipliée par $\mathcal A$ et $B$, cas particulier. La factorisation est démontrée pour un opérateur linéaire abstrait ; le laplacien de \texttt{lap} n'est linéaire que sur les fonctions dérivables, et ce passage n'est pas formalisé.
\end{remarque}

Ne sont pas démontrés ici : que les quatre fonctions $\cosh$, $\sinh$, $\cos$, $\sin$ engendrent toutes les solutions de l'équation du quatrième ordre (les énoncés portent sur cette famille, qui est celle d'Euler) ; les valeurs numériques des racines ($4{,}7300$, $3{,}9266$, $1{,}8751$, $3{,}7502$, $7{,}8532$, \dots) et les intervalles musicaux qu'Euler en tire ; le déterminant $8 \times 8$ des conditions du § 47 et l'existence d'une infinité de modes pour toute position du stylet ; la limite $\lambda \to 0$ ; les §§ 10 à 25 des « Remarques », dont l'équation des bouts encastrés du § 12 et l'écart de sa forme à l'équation exacte hors de $\lambda = \frac12$ ; l'équivalence $c = \frac34 AB$ du § 46 et les anneaux des §§ 53 à 55 ; les séries (E) à (I) et la plaque carrée libre du § 4 ; les feuillets sur la chaleur.

\section[Français 9115 : la courbure des surfaces]{Français 9115 : la courbure des surfaces\protect\verifie{Fr9115Curvature}}\label{sec:fr-9115-courbure}

Le volume Français 9115 s'ouvre sur l'« Exposition des principes qui servent de base à la théorie des surfaces élastiques » (vues 3 à 41) et le « Mémoire sur la courbure des surfaces » (vues 43 à 89), suivis de brouillons sur les « moyennes courbures » (vues 91 à 114) et de feuillets épars, dont un calcul sur la sphère relié parmi les calculs de la chaleur (vue 650). Sophie Germain y propose de mesurer la courbure d'une surface par la somme des inverses des rayons principaux, et son irrégularité par leur différence. On démontre ici les énoncés de la lecture modernisée sur la courbure normale.

\noindent\emph{Conventions.} Pour une surface dont le plan tangent est horizontal au point, $z = \frac12(rx^2 + 2sxy + ty^2) + \dots$, la courbure de la section normale de direction $A$ est $k(A) = r\cos^2A + 2s\sin A\cos A + t\sin^2A$ (\texttt{kq}). En axes principaux, c'est la formule d'Euler $\kappa(\varphi) = \kappa_1\cos^2\varphi + \kappa_2\sin^2\varphi$ (\texttt{kappa}), $\varphi$ étant compté depuis la direction de courbure $\kappa_1$, et $H = \frac12(\kappa_1 + \kappa_2)$ est la courbure moyenne.

\subsection*{Le théorème d'Euler et la loi de répartition}

\begin{enonce}[Euler, Recherches sur la courbure des surfaces, 1760, cité aux vues 16 et 48]
Pour tous $r$, $s$, $t$, il existe $\theta_0$ et $\kappa_2 \leq \kappa_1$ avec $\kappa_1 + \kappa_2 = r + t$, $\kappa_1\kappa_2 = rt - s^2$ et $k(A) = \kappa_1\cos^2(A - \theta_0) + \kappa_2\sin^2(A - \theta_0)$ pour tout $A$.
\end{enonce}

\begin{proof}
Soient $R \geq 0$ et $\phi$ le module et l'argument du nombre complexe $\frac{r-t}2 + is$, de sorte que $R\cos\phi = \frac{r-t}2$ et $R\sin\phi = s$. On prend $\theta_0 = \phi/2$ et $\kappa_{1,2} = \frac{r+t}2 \pm R$. Alors $\kappa_1\kappa_2 = \bigl(\frac{r+t}2\bigr)^2 - R^2$ et $R^2 = \bigl(\frac{r-t}2\bigr)^2 + s^2$. Par duplication, $k(A) = \frac{r+t}2 + \frac{r-t}2\cos 2A + s\sin 2A = \frac{r+t}2 + R\cos(2A - \phi)$, et le second membre vaut $\frac{r+t}2 + R\cos 2(A - \theta_0)$ (\texttt{euler\_principal}).
\end{proof}

\begin{enonce}[Sophie Germain, Exposition, N\textsuperscript{o} 4, \vue{fr-9115}{1}{17} à \vue{fr-9115}{1}{19}]
$k(A) - k(A + \frac\pi2) = (r - t)\cos 2A + 2s\sin 2A$ et $k(\frac\pi2 - A) - k(\pi - A) = -(r - t)\cos 2A + 2s\sin 2A$ ; la somme des deux écarts, $4s\sin 2A$, est nulle en axes principaux.
\end{enonce}

\begin{proof}
On remplace $\cos(A + \frac\pi2)$ par $-\sin A$, $\sin(A + \frac\pi2)$ par $\cos A$, et de même pour les autres angles, puis $\cos 2A$ et $\sin 2A$ par leurs formules de duplication (\texttt{kq\_perp\_diff}, \texttt{kq\_complement\_diff}).
\end{proof}

\begin{enonce}[Sophie Germain, mémoire sur la courbure, N\textsuperscript{o} 3, \vue{fr-9115}{3}{57} à \vue{fr-9115}{4}{62}]
Deux sections normales perpendiculaires ont la somme de courbures $\kappa_1 + \kappa_2$. On a $\kappa(\varphi) = H + \frac{\kappa_1 - \kappa_2}2\cos 2\varphi$, avec $\bigl(\frac{\kappa_1 - \kappa_2}2\bigr)^2 = H^2 - \kappa_1\kappa_2$. La section à $45^\circ$ des plans principaux a la courbure moyenne, et $\kappa(\frac\pi4 + u) + \kappa(\frac\pi4 - u) = \kappa_1 + \kappa_2$ pour tout $u$. Enfin $\kappa(-\varphi) = \kappa(\pi - \varphi) = \kappa(\varphi)$.
\end{enonce}

\begin{proof}
$\cos(\varphi + \frac\pi2)^2 = \sin^2\varphi$ et $\sin(\varphi + \frac\pi2)^2 = \cos^2\varphi$ donnent la somme (\texttt{kappa\_add\_perp}). La forme harmonique vient de $\cos 2\varphi = 2\cos^2\varphi - 1$ et $\sin^2 + \cos^2 = 1$ (\texttt{kappa\_harmonic}, \texttt{amplitude\_sq}). En $\varphi = \frac\pi4$, $\cos 2\varphi = 0$ (\texttt{kappa\_pi\_div\_four}) ; en $\frac\pi4 \pm u$, les cosinus valent $\mp\sin 2u$ et se compensent (\texttt{kappa\_quarter\_symm}). La dernière égalité vient de la parité des carrés (\texttt{kappa\_symm}).
\end{proof}

\begin{remarque}
La formule du N\textsuperscript{o} 2 (\vue{fr-9115}{3}{54}) donne $r = f$, le plus grand rayon, pour $\varphi = 0$ : son angle est compté depuis le plan de moindre courbure, comme la lecture l'a vu. Pour $\varphi' = 45^\circ$ dans la troisième proposition, la direction obtenue est à $90^\circ$ du point de départ et sa courbure est $1/g$, non $1/f$ : c'est $\kappa(\frac\pi2) = \kappa_2$ dans la forme harmonique. La note de l'Exposition (\vue{fr-9115}{1}{19}) sur la sphère de rayon égal au rayon de la section à $45^\circ$ est l'égalité $\kappa(\frac\pi4) = H$.
\end{remarque}

\begin{enonce}[Énoncé de l'édition, d'après les « Idées à vérifier », \vue{fr-9115}{7}{137}]
Si $\kappa_1 \neq \kappa_2$, la section de direction $\varphi$ a la courbure moyenne si et seulement si $\cos 2\varphi = 0$, c'est-à-dire pour les deux sections à $45^\circ$ des plans principaux.
\end{enonce}

\begin{proof}
Par la forme harmonique, $\kappa(\varphi) = H$ équivaut à $\frac{\kappa_1 - \kappa_2}2\cos 2\varphi = 0$ (\texttt{kappa\_eq\_mean\_iff}).
\end{proof}

\subsection*{La moyenne sur toutes les directions}

\begin{enonce}[Sophie Germain, mémoire sur la courbure, N\textsuperscript{os} 4 à 6, \vue{fr-9115}{4}{64} à \vue{fr-9115}{4}{72}]
$\int_0^{2\pi}\kappa(\varphi)\,d\varphi = 2\pi H$ : la moyenne de la courbure normale sur toutes les directions est $H$. La surface des distances, de lignes de distance $\frac12\kappa(\varphi)s^2$ sur le cercle de rayon $s$, a l'étendue $\int_0^{2\pi}\frac12\kappa(\varphi)s^2\cdot s\,d\varphi = 2\pi s\cdot\frac12 Hs^2$ du cylindre de hauteur $\frac12 Hs^2$, la surface des distances moyennes. Enfin, pour un arc d'angle $x$, $|x| \leq 1$, sur un cercle de rayon $\rho \geq 0$, la ligne de distance $\rho(1 - \cos x)$ diffère de $\rho x^2/2$ d'au plus $\frac5{96}\rho x^4$.
\end{enonce}

\begin{proof}
$\int_0^{2\pi}\cos^2 = \int_0^{2\pi}\sin^2 = \pi$, d'où $\pi(\kappa_1 + \kappa_2)$ (\texttt{integral\_kappa}) ; l'intégrande de la surface des distances est $\frac{s^3}2\kappa(\varphi)$ (\texttt{distance\_surface}). Pour la dernière, $\rho(1 - \cos x) - \rho x^2/2 = -\rho\bigl(\cos x - (1 - x^2/2)\bigr)$, et $|\cos x - (1 - x^2/2)| \leq \frac5{96}|x|^4$ pour $|x| \leq 1$ (\texttt{sinus\_verse}).
\end{proof}

\begin{remarque}
Avec $s = \rho x$, $\rho x^2/2 = s^2/(2\rho)$ : à arc égal et infiniment petit, la ligne de distance est proportionnelle à la courbure, ce que dit le N\textsuperscript{o} 4. L'égalité des deux étendues ne demande donc que la moyenne des courbures ; la somme sur « toutes les sections, répétée un nombre infini de fois » de l'Exposition n'a pas de sens, et l'intégrale la remplace. Le développement des deux surfaces sur un plan et les triangles de la vue 75 ne sont pas formalisés.
\end{remarque}

\subsection*{Trois sections normales voisines}

\begin{enonce}[Sophie Germain, feuillet « 55 bis », \vue{fr-9115}{6}{114}]
\begin{align*}
\kappa(A) - \kappa(A+a) &= (\kappa_1-\kappa_2)\sin a\,\sin(2A+a), \\
\kappa(A-a) - \kappa(A) &= (\kappa_1-\kappa_2)\sin a\,\sin(2A-a), \\
\kappa(A-a) + \kappa(A+a) - 2\kappa(A) &= -2(\kappa_1-\kappa_2)\sin^2 a\,\cos 2A .
\end{align*}
Si $\kappa_2 < \kappa_1$, $0 \leq A < \frac\pi4$ et $\sin a \neq 0$, la diminution de $A - a$ à $A$ est plus petite que celle de $A$ à $A + a$. La dérivée $\kappa'(\theta) = -(\kappa_1 - \kappa_2)\sin 2\theta$ vérifie $|\kappa'(\theta)| \leq |\kappa'(\frac\pi4)|$.
\end{enonce}

\begin{proof}
Dans la forme harmonique, on écrit $2A = (2A + a) - a$ et $2(A+a) = (2A + a) + a$ ; la différence $\cos(P - a) - \cos(P + a)$ vaut $2\sin P\sin a$ (\texttt{three\_sections\_one}, et de même \texttt{three\_sections\_two} avec $P = 2A - a$). Pour la troisième, $\cos(2A - 2a) + \cos(2A + 2a) - 2\cos 2A = 2\cos 2A(\cos 2a - 1)$ et $\cos 2a = 1 - 2\sin^2a$ (\texttt{three\_sections\_three}). Dans les conditions énoncées, $\cos 2A > 0$ et $\sin^2a > 0$, et la différence seconde est négative (\texttt{decrease\_compare}). La dérivée se calcule terme à terme, et $|\sin 2\theta| \leq 1 = \sin\frac\pi2$ (\texttt{hasDerivAt\_kappa}, \texttt{abs\_deriv\_kappa\_le}).
\end{proof}

\begin{contrexemple}[à la conclusion du feuillet sur les rayons]
Si $0 < \kappa_2 < \kappa_1$, les rayons des deux sections à $45^\circ$ n'ont pas pour somme celle des rayons principaux : $\frac1{\kappa(\pi/4)} + \frac1{\kappa(\pi/4)} \neq \frac1{\kappa_1} + \frac1{\kappa_2}$.
\end{contrexemple}

\begin{proof}
Le premier membre vaut $4/(\kappa_1 + \kappa_2)$, le second $(\kappa_1 + \kappa_2)/\kappa_1\kappa_2$, et leur égalité donnerait $(\kappa_1 + \kappa_2)(\kappa_1 - \kappa_2)^2 = 0$ (\texttt{radii\_counterexample}).
\end{proof}

\begin{remarque}
L'énoncé juste porte sur les courbures, $\kappa(\frac\pi4 - a) + \kappa(\frac\pi4 + a) = \kappa_1 + \kappa_2$ (\texttt{kappa\_quarter\_symm}). La borne sur la dérivée confirme la dernière phrase du feuillet, que la courbure varie d'autant plus vite qu'on approche de la section de courbure moyenne, et contredit la phrase intermédiaire selon laquelle, à partir de $45^\circ$, l'augmentation serait plus rapide près du plan de plus grande courbure.
\end{remarque}

\subsection*{La somme des courbures sur la sphère}

\begin{enonce}[Sophie Germain, « Somme des raisons inverses des deux rayons de principale courbure », \vue{fr-9115}{33}{650}]
Soit $z = \sqrt{\rho^2 - x^2 - y^2}$ avec $\rho > 0$ et $\rho^2 - x^2 - y^2 > 0$. Ses dérivées partielles sont $p = -x/z$, $q = -y/z$, $r = -(\rho^2 - y^2)/z^3$, $s = -xy/z^3$, $t = -(\rho^2 - x^2)/z^3$, et
\[
\frac{(1+q^2)\,r - 2pq\,s + (1+p^2)\,t}{(1+p^2+q^2)^{3/2}} = -\frac2\rho .
\]
\end{enonce}

\begin{proof}
La dérivée de $\sqrt{D}$, $D = \rho^2 - x^2 - y^2$, est $-2x/(2\sqrt D)$ (\texttt{sphere\_p}, \texttt{sphere\_q}) ; celles de $-x/z$ et de $-y/z$ se calculent par la règle du quotient, en se servant de $z^2 = D$ (\texttt{sphere\_r}, \texttt{sphere\_s}, \texttt{sphere\_t}). Posons $w = 1/z$, de sorte que $Dw^2 = 1$. Alors $1 + p^2 + q^2 = (\rho w)^2$, dont la racine est $\rho w$. Le numérateur vaut $-2\rho^2w^3$, car sa différence avec $-2\rho^2w^3$ est $-w^3(x^2+y^2)(Dw^2 - 1)$. Le quotient est $-2\rho^2w^3/(\rho w)^3 = -2/\rho$ (\texttt{sphere\_sum}).
\end{proof}

\begin{remarque}
La valeur $-2/\rho$ est la somme $\frac1R + \frac1{R'}$ attendue, avec le signe de la normale dirigée vers le bas. Les dérivées secondes sont celles des expressions de $p$ et de $q$ sur le disque ouvert, et la démonstration établit que la formule du feuillet est juste sur la sphère. Le changement de variables qu'elle fait ensuite, en $\theta$ et $\varphi$ à $r$ constant, n'est pas formalisé ; la lecture le juge faux.
\end{remarque}

Ne sont pas démontrés ici : le théorème de Meusnier et ce que l'Exposition dit des sections obliques ; le développement sur un plan des surfaces des distances (vues 72 à 77) ; les surfaces dont les rayons principaux sont de sens contraires (N\textsuperscript{o} 8) ; la comparaison avec l'anneau de Fourier ; la courbure comme quantité dynamique et les voûtes (N\textsuperscript{o} 9) ; le calcul de variation de la vue 483 ; les « Idées à vérifier » sur la chaînette, dont seule la remarque sur les sections à $45^\circ$ est démontrée ; les deux théorèmes proposés aux géomètres (vue 465), qui sont des définitions.

\section[Français 9118 : l'arithmétique des brouillons à Gauss, les notes de lecture, le levier, le centre de gravité, la lame]{Français 9118 : l'arithmétique des brouillons à Gauss, les notes de lecture, le levier, le centre de gravité, la lame\protect\verifie{Fr9118}}\label{sec:fr-9118}

Français 9118 est le dossier extérieur d'une carrière : des lettres reçues, trois brouillons de sa main, des copies, des feuillets détachés. La lecture modernisée y relève plusieurs énoncés mathématiques précis. Les deux brouillons de lettres à Gauss, signés \q{Le Blanc} (\vue{fr-9118}{3}{43} à \vue{fr-9118}{3}{46}), portent une généralisation du théorème de l'article~357 des \emph{Disquisitiones}, le passage de $Y, Z$ à $Y', Z'$, une annonce sur l'équation de Fermat pour $n = p - 1$ et une conjecture sur le déterminant de l'adjointe. Trois feuillets de notes de lecture (\vue{fr-9118}{5}{91}, \vue{fr-9118}{5}{92}, \vue{fr-9118}{5}{93}) traitent de $p^4 + 4$, de $p^4 + q^4$ et d'une composition de formes de Gauss. S'y ajoutent un cours d'une autre main sur le levier et l'infini (\vue{fr-9118}{2}{40} à \vue{fr-9118}{3}{42}), deux rédactions de sa main sur le centre de gravité du triangle (\vue{fr-9118}{5}{91}, \vue{fr-9118}{5}{98}, \vue{fr-9118}{5}{99}), la note de Legendre \q{Pag.~155} sur l'équation d'Euler pour la lame (\vue{fr-9118}{4}{63}, \vue{fr-9118}{4}{64}) et sa lettre du 4~décembre 1811 (\vue{fr-9118}{3}{53}).

\noindent\emph{Conventions.} On note $G_m = \sum_{i<m} X^i = (X^m - 1)/(X - 1) \in \mathbf Z[X]$ ; pour $p$ premier, $G_p$ est le polynôme cyclotomique $\Phi_p$ (\texttt{G\_prime}). Une forme binaire est un triplet $(a, b, c)$ d'entiers, pour $ax^2 + bxy + cy^2$ ; le $(a, b, c)$ de Gauss, $ax^2 + 2bxy + cy^2$, est donc ici $(a, 2b, c)$. Deux formes sont proprement équivalentes si l'une se déduit de l'autre par une substitution entière de déterminant~$1$, équivalentes si le déterminant vaut $\pm 1$.

\subsection*{L'article 357 et la généralisation de Sophie Germain}

\begin{enonce}[Énoncé de l'édition, d'après la vue 43, l'article 357 pour $p = 3, 5, 7$]
$4G_3 = (2X+1)^2 + 3$, $4G_5 = (2X^2+X+2)^2 - 5X^2$ et $4G_7 = (2X^3+X^2-X-2)^2 + 7(X^2+X)^2$.
\end{enonce}

Ce sont des identités de polynômes, vérifiées par le calcul (\texttt{art357\_three}, \texttt{art357\_five}, \texttt{art357\_seven}).

\begin{enonce}[Sophie Germain, vue 43, la généralisation de l'article 357\pisteresolue{fr-9118-gauss-drafts-printed}]
Soient $n \geq 1$ et $c$ un entier impair. Si $4G_n = Y^2 - cZ^2$ avec $Y, Z \in \mathbf Z[X]$, alors pour tout $s \geq 0$ il existe $Y_s, Z_s \in \mathbf Z[X]$ tels que
\[ 4\,\frac{X^{n^s} - 1}{X - 1} = Y_s^2 - c\,Z_s^2 . \]
En particulier $4G_{3^s} = Y^2 + 3Z^2$, $4G_{5^s} = Y^2 - 5Z^2$ et $4G_{7^s} = Y^2 + 7Z^2$ ont des solutions dans $\mathbf Z[X]$ pour tout $s$.
\end{enonce}

\begin{proof}
On suit la démonstration de la lecture. D'abord $G_{ab} = G_a \cdot G_b(X^a)$ : multipliés par $X - 1$, les deux membres valent $X^{ab} - 1$, et $X - 1$ est simplifiable dans $\mathbf Z[X]$ (\texttt{G\_mul}). D'où $G_{n^{s+1}} = G_{n^s} \cdot G_n(X^{n^s})$, et l'hypothèse, transportée par $X \mapsto X^{n^s}$, donne $4G_n(X^{n^s}) = Y(X^{n^s})^2 - cZ(X^{n^s})^2$. Il reste à composer.

Parité (\texttt{parity}) : si $4A = Y^2 - cZ^2$ avec $c = 2k+1$, on réduit modulo~$2$ ; dans $\mathbf F_2[X]$, $0 = \bar Y^2 - \bar Z^2 = (\bar Y - \bar Z)^2$, et comme $\mathbf F_2[X]$ est intègre, $\bar Y = \bar Z$. Le noyau de la réduction $\mathbf Z[X] \to \mathbf F_2[X]$ est l'idéal engendré par $2$ : on a $Y = Z + 2W$.

Composition (\texttt{compose}) : si $4A = Y_1^2 - cZ_1^2$ et $4B = Y_2^2 - cZ_2^2$, écrivons $Y_i = Z_i + 2W_i$. L'identité $(Y_1^2 - cZ_1^2)(Y_2^2 - cZ_2^2) = (Y_1Y_2 + cZ_1Z_2)^2 - c(Y_1Z_2 + Z_1Y_2)^2$ donne $16AB$, et
\[ Y_1Y_2 + cZ_1Z_2 = 2U, \quad U = (k+1)Z_1Z_2 + W_1Z_2 + Z_1W_2 + 2W_1W_2, \qquad Y_1Z_2 + Z_1Y_2 = 2V, \quad V = Z_1Z_2 + W_1Z_2 + Z_1W_2 . \]
Donc $16AB = 4(U^2 - cV^2)$, et $4AB = U^2 - cV^2$ après simplification par~$4$. Une récurrence sur $s$, à partir de $4G_1 = 2^2 - c \cdot 0^2$, achève la preuve (\texttt{generalisation}). Les trois cas particuliers en découlent avec $c = -3, 5, -7$ (\texttt{generalisation\_three}, \texttt{\_five}, \texttt{\_seven}).
\end{proof}

\begin{remarque}
La preuve n'utilise ni que $n$ soit premier ni que $c$ vaille $\pm n$ : il suffit que $c$ soit impair et que la représentation de $4G_n$ soit donnée. La primalité ne sert qu'à l'article~357 lui-même, qui n'est démontré ici que pour $p = 3, 5, 7$. Le signe est bien celui de $n$ pour toutes les valeurs de~$s$, comme le dit la lecture, puisque $c$ ne change pas.
\end{remarque}

\begin{enonce}[Sophie Germain, vue 44, le passage de $Y, Z$ à $Y', Z'$]
Soit $n$ impair. Si $4G_n = Y^2 - cZ^2$, alors $(X+1)\sum_{i<n}(-X)^i = X^n + 1$ et
\[ 4\,\frac{X^n + 1}{X + 1} = Y(-X)^2 - c\,Z(-X)^2 ; \]
le coefficient de degré $i$ de $Y(-X)$ est $(-1)^i$ fois celui de $Y$.
\end{enonce}

On applique à l'hypothèse la substitution $X \mapsto -X$, qui envoie $G_n$ sur $\sum_{i<n}(-X)^i$ (\texttt{passage\_neg}) ; la formule des coefficients se vérifie sur les monômes (\texttt{coeff\_comp\_neg\_X}). Cela revient bien à changer le signe des termes de degré impair.

\begin{remarque}
L'hypothèse \q{$n$ impair}, absente du feuillet et rétablie par la lecture, est nécessaire : pour $n$ pair, $X + 1$ ne divise pas $X^n + 1$, puisque $X^n + 1$ vaut $2$ en $-1$ (\texttt{not\_dvd\_of\_even}).
\end{remarque}

\subsection*{Fermat pour l'exposant $n = p - 1$}

\begin{enonce}[Énoncé de l'édition, d'après la vue 44, le premier cas]
Soit $p$ premier, et $x, y, z$ des entiers dont aucun n'est divisible par $p$. Alors $x^{p-1} + y^{p-1} \neq z^{p-1}$.
\end{enonce}

\begin{proof}
Modulo $p$, le petit théorème de Fermat donne $x^{p-1} \equiv y^{p-1} \equiv z^{p-1} \equiv 1$, et l'égalité donnerait $1 + 1 \equiv 1$, soit $1 \equiv 0$ (\texttt{fermat\_first\_case}).
\end{proof}

La condition $p \equiv 7 \pmod 8$ n'y sert pas, comme le dit la lecture. Les autres affirmations arithmétiques du passage sont vérifiées aussi. Si $p \equiv 7 \pmod 8$, alors $n = p - 1 \equiv 6 \pmod 8$, $4 \nmid n$, $m = (p-1)/2 \equiv 3 \pmod 4$ et $p = 2m + 1$ (\texttt{mod\_eight\_seven}). Pour $p$ premier impair, $p \equiv 7 \pmod 8$ équivaut à ce que $2$ soit un carré et $-1$ ne le soit pas modulo $p$ (\texttt{seven\_mod\_eight\_iff}, par la loi complémentaire de mathlib). Si $3 \mid p - 1$, l'impossibilité pour l'exposant $p - 1$ suit du cas $n = 3$ (\texttt{flt\_of\_three\_dvd}, par le théorème de Fermat pour $n = 3$ de mathlib). Enfin les nombres premiers $p \equiv 7 \pmod 8$ inférieurs à $48$ sont $7, 23, 31, 47$, et $3 \mid p - 1$ pour $7$ et $31$ seulement (\texttt{first\_primes\_seven\_mod\_eight}) : si l'on ne compte comme acquis que les exposants $3$ et $4$, l'annonce n'a de contenu propre qu'à partir de $p = 23$, puis $p = 47$, comme l'écrit la lecture.

\subsection*{Le déterminant de l'adjointe}

\begin{enonce}[Sophie Germain, vue 46, la déterminante de l'adjointe]
Pour toute matrice $A$ carrée d'ordre $n$ sur un anneau commutatif, $\det(\operatorname{adj} A) = (\det A)^{n-1}$. Dans la convention de Gauss pour les formes ternaires, où la déterminante est $D = -\det M$ et l'adjointe $-\operatorname{adj} M$, la déterminante de l'adjointe est $D^2$.
\end{enonce}

La première assertion est le lemme \texttt{Matrix.det\_adjugate} de mathlib, récrit pour $\mathrm{Fin}\,n$ (\texttt{det\_adjugate\_fin}) ; la seconde en suit, puisque $-\det(-\operatorname{adj} M) = -(-1)^3 (\det M)^2$ (\texttt{gauss\_ternary\_adjoint}). Pour $n = 2$, la formule donne $D^{n-1} = D$, ce que la lecture retient.

\begin{enonce}[Énoncé de l'édition, d'après la vue 46, le radical]
Une matrice symétrique $M$ d'ordre $n$ sur un corps, de déterminant nul, admet un vecteur $v \neq 0$ tel que $Q(x + tv) = Q(x)$ pour tout $x$ et tout $t$, où $Q(x) = {}^tx\,M x$.
\end{enonce}

On prend $v$ non nul avec $Mv = 0$ ; par symétrie ${}^tv\,M x = {}^t(Mv)\,x = 0$, et les termes croisés de $Q(x + tv)$ s'annulent (\texttt{radical\_of\_det\_eq\_zero}). C'est le premier pas de la réduction d'une forme de déterminant nul à une forme en moins de variables ; le changement de base qui achève cette réduction n'est pas formalisé.

\subsection*{$p^4 + 4$}

\begin{enonce}[Sophie Germain, vue 91, $p^4 + 4$\piste{fr-9118-p4-plus-4-identity}]
Pour tout entier naturel $p$, $p^4 + 4$ est premier si et seulement si $p = 1$. On a $p^4 + 4 = (p^2)^2 + 2^2 = (p^2 - 2)^2 + (2p)^2$, et les deux paires $\{p^2, 2\}$ et $\{|p^2 - 2|, 2p\}$ coïncident si et seulement si $p \leq 2$.
\end{enonce}

\begin{proof}
L'identité de Sophie Germain $a^4 + 4b^4 = (a^2 + 2b^2 - 2ab)(a^2 + 2b^2 + 2ab)$ est une identité de polynômes (\texttt{germain\_identity}). Pour $p = q + 1$ elle s'écrit $p^4 + 4 = (q^2 + 1)\big((q+2)^2 + 1\big)$ ; si ce produit est premier, l'un des facteurs vaut $1$, et seul $q^2 + 1$ le peut, d'où $q = 0$ ; et $0^4 + 4 = 4$ n'est pas premier (\texttt{p4\_add\_four\_prime\_iff}). Pour les paires : la première coïncidence demande $2 = 2p$, et la seconde $p^2 = 2p$, soit $p \in \{0, 2\}$ ; réciproquement $p = 0, 1, 2$ conviennent (\texttt{same\_pair\_iff}).
\end{proof}

\begin{remarque}
La lecture relève que les deux paires coïncident pour $p = 1$ et $p = 2$. Elles coïncident aussi pour $p = 0$ ($4 = 0 + 4 = 4 + 0$), cas que le feuillet écarte sans doute d'emblée. La factorisation n'est pas sur le feuillet, qui procède par les deux sommes de deux carrés et le théorème d'Euler sur les deux représentations ; ce théorème n'est pas formalisé ici, et la factorisation règle la question sans lui.
\end{remarque}

\subsection*{$p^4 + q^4$ et les nombres de Fermat}

\begin{enonce}[Sophie Germain, vue 93, $p^4 + q^4$]
On a $p^4 + q^4 = (p^2 - q^2)^2 + 2(pq)^2 = (p^2 + q^2)^2 - 2(pq)^2$ dans tout anneau commutatif. Pour $i \geq 2$, $2^{2^i} + 1 = (2^{2^{i-2}})^4 + 1^4$, tandis que $3$ et $5$ ($i = 0, 1$) ne sont pas sommes de deux puissances quatrièmes d'entiers naturels.
\end{enonce}

Les identités sont vérifiées par le calcul (\texttt{p4\_add\_q4\_identities}). Pour les nombres de Fermat, $2^{2^{j+2}} = 2^{4 \cdot 2^j} = (2^{2^j})^4$ (\texttt{fermat\_number\_eq}) ; pour $3$ et $5$, un entier $p \geq 2$ donne déjà $p^4 \geq 16$, et il reste à essayer $p, q \in \{0, 1\}$ (\texttt{three\_five\_not\_sum\_fourth}). La restriction $i \geq 2$, que pose la lecture, est nécessaire.

\begin{contrexemple}[à la lecture et au feuillet, vue 93, l'unicité pour $x^2 - 2y^2$]
Le nombre premier $17 = 2^4 + 1^4$ s'écrit $5^2 - 2 \cdot 2^2$, qui est $(p^2 + q^2)^2 - 2(pq)^2$, mais aussi $23^2 - 2 \cdot 16^2$. Plus généralement, toute représentation $N = x^2 - 2y^2$ en donne une autre, $N = (3x + 4y)^2 - 2(2x + 3y)^2$.
\end{contrexemple}

Les deux égalités sont vérifiées numériquement (\texttt{x\_sq\_sub\_two\_y\_sq\_not\_unique}), la seconde formule par le calcul (\texttt{pell\_step}) ; elle traduit la multiplication par l'unité $3 + 2\sqrt 2$ de $\mathbf Z[\sqrt 2]$. Le feuillet affirme que, $p^4 + q^4$ étant premier, il n'est \q{que de ces seules manières} un carré plus ou moins un double carré ; la lecture précise que la représentation par $x^2 - 2y^2$ est unique. C'est faux : elle n'est unique qu'à multiplication près par les unités, qui sont en nombre infini. La remarque vaut pour $x^2 + 2y^2$, où les unités sont $\pm 1$. On le vérifie ici pour $17 = 3^2 + 2 \cdot 2^2$ et $257 = 15^2 + 2 \cdot 4^2$, seules représentations en entiers naturels (\texttt{seventeen\_unique}, \texttt{two\_fifty\_seven\_unique}) ; l'énoncé général pour tous les nombres premiers, qui demande que $\mathbf Z[\sqrt{-2}]$ soit principal, n'est pas formalisé.

\subsection*{L'objection à Gauss sur la composition des formes}

La composition est prise ici au sens de Dirichlet. On dit que $F$ est composée de $f$ et de $g$ s'il existe $a_1, a_2, B, c$ avec $\operatorname{pgcd}(a_1, a_2) = 1$ tels que $f$ soit proprement équivalente à $(a_1, B, a_2c)$, $g$ à $(a_2, B, a_1c)$, et $(a_1a_2, B, c)$ à $F$. L'identité
\[ (a_1x^2 + Bxy + a_2c\,y^2)(a_2z^2 + Bzw + a_1c\,w^2) = a_1a_2X^2 + BXY + cY^2, \quad X = xz - cyw,\ Y = a_1xw + a_2yz + Byw \]
justifie le mot (\texttt{dirichlet\_identity}).

\begin{enonce}[Énoncé de l'édition, d'après la vue 92, les quatre compositions\piste{fr-9118-composition-objection}]
En notation de Gauss : $(16, 3, 19) \circ (8, 1, 37) = (8, -1, 37)$ ; $(16, -3, 19) \circ (8, 1, 37) = (13, -2, 23)$ ; $(8, -1, 37) \circ (5, 0, 59) = (13, -2, 23)$ ; et $(5, 0, 59) \circ (5, 0, 59) = (1, 0, 295)$. Les formes $(16, \pm 3, 19)$ sont équivalentes (improprement), mais $(8, -1, 37)$ et $(13, -2, 23)$ ne sont pas équivalentes, ni proprement ni improprement.
\end{enonce}

\begin{proof}
Pour la première, $(16, 6, 19)$ devient $(19, 146, 296)$ par la substitution $(x, y) \mapsto (-y, x + 4y)$, $(8, 2, 37)$ devient $(8, 146, 703)$ par $(x, y) \mapsto (x + 9y, y)$ ; avec $a_1 = 19$, $a_2 = 8$, $B = 146$, $c = 37$, la composée $(152, 146, 37)$ se réduit à $(8, -2, 37)$ par $(x, y) \mapsto (-x, 2x - y)$. Les trois autres se traitent de même, avec $(a_1, a_2, B, c) = (19, 8, 82, 13)$, $(8, 5, 30, 13)$ et $(5, 59, 0, 1)$ ; toutes les substitutions sont de déterminant~$1$ et chaque égalité de formes est une identité vérifiée par le calcul (\texttt{composite\_one} à \texttt{composite\_four}). La substitution $(x, y) \mapsto (x, -y)$, de déterminant $-1$, échange $(16, 3, 19)$ et $(16, -3, 19)$ (\texttt{opposite\_equiv}). Enfin $(8, -1, 37)$ représente $8$ en $(1, 0)$ ; si elle se déduisait de $13x^2 - 4xy + 23y^2$ par une substitution, celle-ci représenterait $8$, soit $104 = (13u - 2v)^2 + 295v^2$, d'où $v = 0$ et $169u^2 = 104$, impossible (\texttt{not\_equiv\_8\_13}).
\end{proof}

\begin{enonce}[Énoncé de l'édition, d'après la vue 92, les formes réduites]
Les formes $(a, b, c)$ de Gauss de déterminant $b^2 - ac = -295$, réduites ($a > 0$, $|2b| \leq a \leq c$, $b \geq 0$ si $|2b| = a$ ou $a = c$) et proprement primitives ($\operatorname{pgcd}(a, 2b, c) = 1$) sont les huit formes $(1, 0, 295)$, $(5, 0, 59)$, $(8, \pm 1, 37)$, $(13, \pm 2, 23)$, $(16, \pm 3, 19)$.
\end{enonce}

\begin{proof}
De $4b^2 \leq a^2$ et $ac \geq a^2$ on tire $295 = ac - b^2 \geq \tfrac34 a^2$, donc $a \leq 19$ et $|b| \leq 9$ ; $c$ est alors déterminé par $ac = b^2 + 295$, et l'on examine les cas restants (\texttt{reduced\_forms\_295}).
\end{proof}

\begin{remarque}
Sans la condition de primitivité propre, on trouve huit formes de plus, $(2, 1, 148)$, $(4, \pm 1, 74)$, $(8, \pm 3, 38)$, $(10, 5, 32)$, $(16, \pm 5, 20)$, qui forment l'ordre improprement primitif ; la lecture parle à juste titre de l'ordre \q{où Gauss compose}. Ce qui est démontré confirme le calcul de la lecture : le produit de Gauss est le bon, et la seconde forme de Sophie Germain provient de la forme opposée du premier facteur. Le contre-exemple \texttt{not\_equiv\_8\_13} montre de plus que les deux résultats ne sont pas même improprement équivalents, de sorte qu'ils ne peuvent être deux valeurs d'un même produit de classes. Ne sont pas démontrés : que la composition de Dirichlet passe aux classes d'équivalence propre et coïncide avec la composition de l'article~235, et que deux formes réduites distinctes ne soient jamais proprement équivalentes, d'où le nombre de classes $8$.
\end{remarque}

\subsection*{Le levier, l'infini et l'hyperbole}

\begin{enonce}[Énoncé de l'édition, d'après les vues 40 à 42, les formes indéterminées]
Pour tout réel $c$, $c/x \to 0$ et $x \to +\infty$ quand $x \to +\infty$, tandis que $(c/x) \cdot x \to c$ ; de même $c/x \to 0$ et $1/x \to 0$, tandis que $(c/x)/(1/x) \to c$.
\end{enonce}

Pour $x > 0$ les deux expressions valent exactement $c$ (\texttt{indeterminate\_product}, \texttt{indeterminate\_quotient}). Chaque hyperbole $xy = c$ réalise ainsi une valeur de \q{$0 \times \infty$}, et \q{$0/0$} vaut \q{tout ce qu'on veut}, ce que le cours dit en refusant l'infini actuel. La première limite est aussi celle du levier : la force qui équilibre un moment donné à la distance $x$ tend vers $0$ quand $x$ croît.

\subsection*{Le centre de gravité du triangle}

\begin{enonce}[Sophie Germain, vues 91, 98 et 99, le centre de gravité du triangle]
L'équation $2x = (a + x)/2$ équivaut à $x = a/3$. Dans un espace vectoriel réel, soient $D$, $E$, $F$ les milieux de $BC$, $CA$, $AB$, et $g(P, Q, R) = \tfrac13(P + Q + R)$. Alors
\[ \tfrac14\, g(B, F, D) + \tfrac14\, g(C, E, D) + \tfrac12 \cdot \tfrac12 (A + D) = g(A, B, C), \qquad g(A, B, C) = D + \tfrac13 (A - D), \]
et $g(B, F, D) = B + \tfrac12\big(g(A, B, C) - B\big)$.
\end{enonce}

Toutes ces égalités sont linéaires et se vérifient par le calcul (\texttt{centroid\_equation}, \texttt{centroid\_decomposition}, \texttt{centroid\_on\_median}, \texttt{small\_centroid\_homothety}). La première égalité vectorielle est la décomposition du feuillet : le parallélogramme $AFDE$, de centre le milieu de $AD$ et de poids la moitié, et les deux petits triangles, d'un quart chacun ; la deuxième dit que le centre de gravité est au tiers de la médiane à partir de la base ; la troisième est la similitude de rapport $\tfrac12$ que le feuillet invoque pour placer les centres des petits triangles. Les poids, c'est-à-dire les aires, sont posés et non calculés.

\subsection*{Legendre, \q{Pag.~155} : la racine parasite et le cas $\lambda = 1/3$}

On note $F_\lambda(\omega) = 2 - 2e^{2\omega} - (e^{\lambda\omega} - e^{(2-\lambda)\omega})(e^{\lambda\omega} - e^{-\lambda\omega})(\cot\lambda\omega + \cot(1-\lambda)\omega)$ le second membre de l'équation d'Euler telle que Legendre la récrit, et
\[ H(\omega) = \sin\tfrac{\omega}{2}\,(2 - 2e^{2\omega}) - 2\,(e^{\omega/2} - e^{3\omega/2})(e^{\omega/2} - e^{-\omega/2})\cos\tfrac{\omega}{2} . \]

\begin{enonce}[Legendre, vues 63 et 64, la racine parasite $\sin\frac12\omega = 0$]
Là où $\sin\frac{\omega}{2} \neq 0$, on a $F_{1/2}(\omega) = H(\omega)/\sin\frac{\omega}{2}$. Pour $\omega = 2k\pi$, $k$ entier non nul, $\sin\frac{\omega}{2} = 0$ mais $H(\omega) \neq 0$.
\end{enonce}

\begin{proof}
La première égalité s'obtient en écrivant $\cot = \cos/\sin$ et en réduisant au même dénominateur (\texttt{parasite\_root}). En $\omega = 2k\pi$, le premier terme de $H$ s'annule et il reste $-2(e^{k\pi} - e^{3k\pi})(e^{k\pi} - e^{-k\pi})\cos k\pi$ ; les deux différences sont non nulles parce que l'exponentielle est injective et $k\pi \neq 0$, et $\cos k\pi \neq 0$ parce que $\sin k\pi = 0$ (\texttt{parasite\_H\_ne\_zero}).
\end{proof}

\begin{remarque}
L'équation multipliée $\sin\frac{\omega}{2}\,H(\omega) = 0$ admet donc les racines $2k\pi$, qui viennent du seul facteur introduit par la multiplication : c'est le diagnostic de Legendre. Que $|F_{1/2}|$ tende vers l'infini au voisinage de ces points, comme le dit la lecture, en résulte par continuité de $H$, mais la limite n'est pas formalisée, non plus que le cas $\omega \to 0$.
\end{remarque}

\begin{enonce}[Legendre, vue 64, la réduction pour $\omega$ grand]
Pour tous $\lambda$ et $\omega$, avec $\varepsilon(\omega) = e^{(2\lambda - 2)\omega} - e^{-2\omega} + e^{-2\lambda\omega}$ et $S = \cot\lambda\omega + \cot(1-\lambda)\omega$,
\[ e^{-2\omega}\,F_\lambda(\omega) = 2e^{-2\omega} - 2 + \big(1 - \varepsilon(\omega)\big)\,S , \]
et, si $0 < \lambda < 1$, $\varepsilon(\omega) \to 0$ quand $\omega \to +\infty$.
\end{enonce}

Le produit des deux premiers facteurs, multiplié par $e^{-2\omega}$, vaut exactement $\varepsilon(\omega) - 1$ (\texttt{euler\_large\_omega}) ; les trois exponentielles de $\varepsilon$ ont des exposants négatifs (\texttt{eps\_tendsto\_zero}). L'équation $F_\lambda = 0$ s'écrit donc $(1 - \varepsilon)S = 2 - 2e^{-2\omega}$, qui tend vers $S = 2$ : c'est la réduction de Legendre, avec son terme d'erreur explicite.

\begin{enonce}[Legendre, vue 64, le cas $\lambda = 1/3$]
Soit $\theta$ avec $\sin 2\theta \neq 0$, et $x = \cot\theta$. Alors $\cot\theta + \cot 2\theta = 2$ équivaut à $3x^2 - 4x - 1 = 0$, qui équivaut à $x = (2 \pm \sqrt 7)/3$.
\end{enonce}

\begin{proof}
On a $\cos 2\theta = \cos^2\theta - \sin^2\theta$ et $\sin 2\theta = 2\sin\theta\cos\theta$, d'où $\cot 2\theta = (x^2 - 1)/(2x)$, avec $x \neq 0$ puisque $\cos\theta \neq 0$. L'équation devient $2x^2 + x^2 - 1 = 4x$. Enfin $3x^2 - 4x - 1 = 3\big(x - \tfrac{2+\sqrt7}{3}\big)\big(x - \tfrac{2-\sqrt7}{3}\big)$ (\texttt{lambda\_third}).
\end{proof}

\begin{remarque}
Avec $\theta = \omega/3$, c'est le calcul de la note, qui est juste. Les valeurs numériques que donne la lecture en note ne sont pas vérifiées en Lean ; un calcul direct donne $\alpha \approx 0{,}5734$ et $\beta \approx 1{,}7828$, donc $3\beta \approx 5{,}348$, et non $\beta \approx 1{,}781$ : la lecture arrondit mal la seconde racine, sans conséquence pour le raisonnement.
\end{remarque}

\subsection*{La vérification par la lame}

\begin{enonce}[Énoncé de l'édition, d'après la vue 53, la condition $dz/dy = 0$]
Soit $z(x, y)$ dérivable en $y$ pour tout $x$. Si $\partial z/\partial y$ est identiquement nulle, alors $z(x, y) = z(x, 0)$ pour tous $x, y$.
\end{enonce}

C'est le théorème des accroissements finis, sous la forme \texttt{is\_const\_of\_deriv\_eq\_zero} de mathlib (\texttt{indep\_of\_deriv\_zero}).

\begin{remarque}
La lecture écrit en note que la condition de Legendre \q{ne suffit pas telle quelle} parce qu'elle n'annule que la dérivée première. Prise comme une identité, ce qui est le sens d'une lame de largeur constante, elle suffit : $z$ ne dépend alors pas de $y$, et toutes les dérivées en $y$ s'annulent. La réserve de la lecture ne vaut que si la condition est lue en un seul point. La réduction de $\Delta^2 z$ à $\partial^4 z/\partial x^4$ qui s'ensuit n'est pas formalisée.
\end{remarque}

Ne sont pas démontrés ici : le théorème de l'article~357 pour $p$ premier quelconque, qui demande les sommes de Gauss ; les deux démonstrations de Sophie Germain, absentes du volume ; l'annonce sur Fermat pour $p \equiv 7 \pmod 8$ dans le second cas ; la réduction de la quantité de Lagrange (1775) à la forme $t^2 - 11u^2$, dont les deux copies du feuillet diffèrent ; la réduction d'une forme de déterminant nul à moins de variables au-delà du radical ; le théorème d'Euler sur les deux représentations comme somme de deux carrés ; l'unicité de la représentation par $x^2 + 2y^2$ pour tout nombre premier ; la théorie des classes de formes (composition sur les classes, unicité de la forme réduite, nombre de classes) ; les limites de la discussion de la racine parasite ; les erreurs de signe d'Euler ($\gamma = \gamma'$) et les lignes nodales, qui demandent le texte d'Euler ; l'équation de la plaque, la courbure moyenne et la note de Libri, que cette section ne traite pas.

\section[NAF 4073 : le triangulaire d'un carré]{NAF 4073 : le triangulaire d'un carré\protect\verifie{Naf4073}}\label{sec:naf-4073}

NAF 4073 est un volume de lettres. La seule ligne de mathématiques qu'il contienne est écrite tête-bêche au haut du feuillet d'adresse de la lettre non signée du 13~mars 1814, d'une main qui n'est pas celle de la lettre (\vue{naf-4073}{3}{42}) : l'égalité
\[ \Big(\frac{n^2 - n}{2}\Big)^2 + \Big(\frac{n^2 + n}{2}\Big)^2 = \frac{n^4 + n^2}{2}, \]
avec la phrase qui la lit, \q{le triangulaire d'un quarré est la somme des quarrés des deux Triangulaires consécutifs}. Sous l'égalité, une ligne isolée que la transcription lit $1 + 64 + 49$. La lecture modernisée note $T_k = k(k+1)/2$ et lit l'égalité comme $T_{n-1}^2 + T_n^2 = T_{n^2}$.

\begin{enonce}[Note de NAF 4073, vue 42, le triangulaire d'un carré\piste{naf-4073-triangular-note}]
Dans tout corps où $2 \neq 0$, $\big(\frac{n^2 - n}{2}\big)^2 + \big(\frac{n^2 + n}{2}\big)^2 = \frac{n^4 + n^2}{2}$. En entiers naturels, pour tout $n$,
\[ T_n + T_{n+1} = (n+1)^2, \qquad T_{n+1} - T_n = n + 1, \qquad T_n^2 + T_{n+1}^2 = T_{(n+1)^2}, \]
et, pour $n \geq 1$, $\binom{n}{2}^2 + \binom{n+1}{2}^2 = \binom{n^2+1}{2}$.
\end{enonce}

\begin{proof}
La première égalité est une identité de fractions rationnelles (\texttt{leaf\_identity}). En entiers, la division par $2$ est exacte puisque $k(k+1)$ est pair, d'où $2T_k = k(k+1)$ (\texttt{two\_mul\_T}) et $T_k = \binom{k+1}{2}$ (\texttt{T\_eq\_choose}). Posons $a = T_n$, $b = T_{n+1}$ : $2a + 2b = n(n+1) + (n+1)(n+2) = 2(n+1)^2$ et $2b - 2a = 2(n+1)$ (\texttt{T\_add\_T}, \texttt{T\_succ\_sub\_T}). L'identité du parallélogramme donne alors $4(a^2 + b^2) = (2a)^2 + (2b)^2 = (n+1)^2\big(n^2 + (n+2)^2\big) = 2(n+1)^2\big((n+1)^2 + 1\big) = 4\,T_{(n+1)^2}$, et l'on simplifie par $4$ (\texttt{triangular\_of\_square}). La forme binomiale en découle par $T_k = \binom{k+1}{2}$, en posant $n = m + 1$ (\texttt{triangular\_of\_square\_choose}).
\end{proof}

\begin{remarque}
L'énoncé est démontré sous la forme $T_n^2 + T_{n+1}^2 = T_{(n+1)^2}$, qui est celle de la lecture avec $n + 1$ au lieu de $n$ ; elle évite la soustraction tronquée des entiers naturels et vaut pour tout $n$, y compris $n = 0$. Aucune hypothèse n'est superflue ni manquante. La preuve ne dit rien de la main ni de la date de la note, qui font l'objet de la piste : celle-ci porte sur l'existence de la note dans ce volume, non sur l'identité, qui est classique. Germain connaissait l'identité : un manuscrit qui lui est attribué, daté par le libraire Quaritch de 1797 environ et reproduit par G.~L.~Alexanderson (\emph{Bulletin of the AMS}, t.~49, 2012, p.~327--331), la donne sous la même forme, après un « J'ai aussi remarqué que ».
\end{remarque}

\begin{enonce}[Énoncé de l'édition, d'après la vue 42, la conséquence par Nicomaque]
Pour tout $n$, $\sum_{i=0}^{n} i^3 = T_n^2$, et $T_{(n+1)^2} + (n+1)^3 = 2\,T_{n+1}^2$, c'est-à-dire $T_{m^2} = 2T_m^2 - m^3$ pour $m \geq 1$.
\end{enonce}

\begin{proof}
Par récurrence : $4\big(T_n^2 + (n+1)^3\big) = n^2(n+1)^2 + 4(n+1)^3 = (n+1)^2(n+2)^2 = 4\,T_{n+1}^2$ (\texttt{nicomachus}). Le même calcul donne $T_{n+1}^2 = T_n^2 + (n+1)^3$, qu'on porte dans l'énoncé précédent (\texttt{triangular\_of\_square\_cubes}).
\end{proof}

\begin{enonce}[Énoncé de l'édition, d'après la vue 42, la ligne $1 + 64 + 49$]
On a $1 + 64 + 49 = 114$ ; $114$ n'est ni un nombre triangulaire ni un carré, et $1 + 64 \neq 49$.
\end{enonce}

Si $T_k = 114$, alors $k(k+1) = 228$, ce qui force $k \leq 15$, et aucun $k \leq 15$ ne convient ; de même $m^2 = 114$ force $m \leq 11$ (\texttt{line\_114}). Le constat de la lecture se trouve confirmé : aucune interprétation de ces trois nombres n'instancie l'identité. Ce que la ligne additionne reste inconnu.

Ne sont pas démontrés ici : l'attribution de la main et la date de la note, qui relèvent du fac-similé et non d'une preuve.


\end{document}
